\newcommand{\sorryleaf}[1]{%
  \par\smallskip\noindent\textcolor{red!68!black}{%
  \textsf{Proof status: statement-only leaf; proof deferred.}}%
  \par\noindent\textcolor{red!68!black}{\textsf{Mathematical source:} }#1\par}

\newcommand{\BlueprintLocalDeclarationRegistry}[4]{%
  \par\smallskip
  {%
  \scriptsize
  \setlength{\tabcolsep}{3pt}%
  \renewcommand{\arraystretch}{1.08}%
  \noindent
  \begin{tabular}{@{}%
    p{0.22\linewidth}%
    p{0.38\linewidth}%
    p{\dimexpr0.40\linewidth-12pt\relax}@{}}
    \toprule
    \textbf{#1} & \textbf{Direct uses} & \textbf{Exact source}\\
    \midrule
    \hyperref[#2]{#1~\ref*{#2}}\newline\hyperref[#2]{\nolinkurl{#2}} &
    \UseLinks{#3} &
    #4\\
    \bottomrule
  \end{tabular}
  }%
  \par\smallskip
}

\newcommand{\sketchslot}[1]{%
  \par\smallskip
  {%
  \footnotesize
  \setlength{\tabcolsep}{4pt}%
  \renewcommand{\arraystretch}{1.05}%
  \noindent
  \begin{tabular}{@{}p{0.24\linewidth}%
    p{\dimexpr0.76\linewidth-8pt\relax}@{}}
    \textcolor{black!65}{\textbf{Future detailed sketch}} &
    \if\relax\detokenize{#1}\relax
      \textcolor{black!45}{\rule{0.94\linewidth}{0.35pt}}
    \else
      #1
    \fi
  \end{tabular}
  }%
  \par\smallskip
}

\newcommand{\Fbar}{\overline{\mathbf F}_q}
\newcommand{\Qlb}{\overline{\mathbf Q}_\ell}
\newcommand{\Zl}{\mathbf Z_\ell}
\newcommand{\RGamma}{R\Gamma}
\newcommand{\RGammac}{R\Gamma_c}
\newcommand{\Sh}{\operatorname{Sh}}
\newcommand{\Kar}{\operatorname{Kar}}
\newcommand{\Wt}{\operatorname{Wt}}

\chapter{Global Sketch}

\section{Arithmetic and geometric conventions}

\begin{definition}[Finite field and geometric fiber]
\label{def:weil-base-field}
Let \(p\) be a prime, let \(a\geq1\), and put \(q=p^a\).  Set
\[
 k_0=\Fq,
 \qquad
 k=\Fbar,
\]
and fix a prime \(\ell\neq p\).  For a finite-type \(k_0\)-scheme \(Y_0\),
write
\[
 Y=Y_0\times_{k_0}k,
 \qquad
 a_Y:Y\longrightarrow\Spec k
\]
for its geometric fiber and structure morphism.
\end{definition}

\BlueprintLocalDeclarationRegistry{Definition}{def:weil-base-field}{none}{Exact local declaration (Definition~\ref{def:weil-base-field}).}

\begin{definition}[Geometric hypotheses]
\label{def:weil-geometric-hypotheses}
A finite-type \(k_0\)-scheme \(Y_0\) is \emph{smooth of pure dimension
\(d\)} if \(Y_0\to\Spec k_0\) is smooth and every irreducible component of
its geometric fiber has Krull dimension \(d\).  It is \emph{proper} if its
structure morphism is separated, of finite type, and universally closed.  It
is \emph{geometrically connected} if \(Y_0\times_{k_0}k'\) is connected for
every field extension \(k'/k_0\), equivalently if \(Y\) is connected.
\uses{def:weil-base-field}
\end{definition}

\BlueprintLocalDeclarationRegistry{Definition}{def:weil-geometric-hypotheses}{def:weil-base-field}{Exact local declaration (Definition~\ref{def:weil-geometric-hypotheses}).}

\begin{definition}[Arithmetic and geometric Frobenius]
\label{def:weil-frobenius-convention}
The arithmetic Frobenius of \(k\) is
\[
 \varphi_q:k\longrightarrow k,
 \qquad z\longmapsto z^q.
\]
The geometric Frobenius element is
\(\gamma=\operatorname{Frob}_q=\varphi_q^{-1}\).  If
\(Y=Y_0\times_{k_0}k\), the associated semilinear automorphism is
\[
 g_Y=\operatorname{id}_{Y_0}\times\Spec(\gamma):Y\longrightarrow Y.
\]
It is distinct from the absolute \(q\)-power morphism of the scheme \(Y\).
Every Frobenius operator below is induced by \(g_Y\) and is geometric.
\uses{def:weil-base-field}
\end{definition}

\BlueprintLocalDeclarationRegistry{Definition}{def:weil-frobenius-convention}{def:weil-base-field}{Lecture 46, Definition 2.1.}

\section{Triangulated coefficient categories}

\begin{convention}[Size]
\label{conv:weil-size}
Fix Grothendieck universes \(\mathcal U\in\mathcal V\) containing the
schemes and coefficient rings under consideration.  Replace every site by a
\(\mathcal V\)-small skeleton within the chosen bound.  Derived categories,
filtered colimits, and Karoubi completions are formed in \(\mathcal V\).
\end{convention}

\BlueprintLocalDeclarationRegistry{Convention}{conv:weil-size}{none}{Exact local convention (Convention~\ref{conv:weil-size}).}

\begin{definition}[Small pro-etale site]
\label{def:weil-small-proetale-site}
For a scheme \(Y\), the site \(Y_{\mathrm{proet}}\) has as objects the
\(\mathcal U\)-bounded weakly etale morphisms \(U\to Y\).  A morphism is
weakly etale when it is flat and its diagonal is flat.  A family
\(\{U_i\to U\}_{i\in I}\) is a covering if every \(U_i\to U\) is weakly
etale and, for every affine open \(V\subset U\), finitely many affine opens
\(V_j\subset U_{i_j}\) have images covering \(V\).

For a ring \(A\), let \(A_Y\) denote the constant sheaf of rings on
\(Y_{\mathrm{proet}}\), and let \(\Sh(Y_{\mathrm{proet}},A)\) be the abelian
category of sheaves of \(A_Y\)-modules.  Its derived category
\(D(Y_{\mathrm{proet}},A)\) is the localization of the category of complexes
at quasi-isomorphisms; it is used only through its ordinary triangulated
structure.
\uses{conv:weil-size}
\end{definition}

\BlueprintLocalDeclarationRegistry{Definition}{def:weil-small-proetale-site}{conv:weil-size}{Lecture 40, Definition 2.3.}

\begin{definition}[Finite-torsion constructibility]
\label{def:weil-finite-torsion-constructible}
Let \(E/\Ql\) be a finite extension, let \(\mathcal O_E\) be its valuation
ring, choose a uniformizer \(\varpi_E\), and put
\[
 A_{E,n}=\mathcal O_E/\varpi_E^n
 \qquad(n\geq1).
\]
A sheaf of \(A_{E,n}\)-modules on \(Y_{\mathrm{proet}}\) is
\emph{constructible} if \(Y\) has a finite partition by constructible locally
closed subschemes on each of which the sheaf is locally constant with finite
\(A_{E,n}\)-module stalks.  The category
\[
 D^b_c(Y,A_{E,n})\subset D(Y_{\mathrm{proet}},A_{E,n})
\]
is the full triangulated subcategory of complexes with finitely many nonzero
cohomology sheaves, all constructible.
\uses{def:weil-small-proetale-site}
\end{definition}

\BlueprintLocalDeclarationRegistry{Definition}{def:weil-finite-torsion-constructible}{def:weil-small-proetale-site}{Lecture 40, Definitions 3.5 and 7.4.}

\begin{definition}[Derived-complete integral coefficients]
\label{def:weil-derived-complete-integral}
The completed coefficient sheaf is
\[
 \widehat{\mathcal O}_{E,Y}=\varprojlim_n(A_{E,n})_Y
\]
in the abelian category of sheaves of rings.  By convention,
\(\Sh(Y_{\mathrm{proet}},\mathcal O_E)\) means modules over
\(\widehat{\mathcal O}_{E,Y}\), and
\(D(Y_{\mathrm{proet}},\mathcal O_E)\) is the derived category of that
abelian category; the notation never means modules over the discrete constant
sheaf \((\mathcal O_E)_Y\).

Let \(\operatorname{Tow}_E(Y)\) be the abelian category of inverse systems
\(C_{n+1}\to C_n\) in
\(\Sh(Y_{\mathrm{proet}},\mathcal O_E)\).  The left exact functor
\[
 \varprojlim:\operatorname{Tow}_E(Y)
 \longrightarrow\Sh(Y_{\mathrm{proet}},\mathcal O_E)
\]
has a total right-derived functor, obtained by a K-injective replacement in
the category of complexes of towers.  It is denoted \(R\varprojlim\).  For a
strict tower of complexes \(C=(C_n)\), its value has the functorial Milnor
triangle
\[
 R\varprojlim_nC_n\longrightarrow
 \prod_{n\geq1}C_n\xrightarrow{1-\mathrm{shift}}
 \prod_{n\geq1}C_n\longrightarrow(R\varprojlim_nC_n)[1].
\]
This triangle is functorial in the strict tower (C).

An object \(K\in D(Y_{\mathrm{proet}},\mathcal O_E)\) is
\emph{derived complete} if
the canonical map
\[
 K\longrightarrow R\varprojlim_n
 \operatorname{Res}_n
 \bigl(K\otimes^L_{\mathcal O_E}A_{E,n}\bigr)
\]
is an isomorphism.  Here
\[
 \operatorname{Res}_n:
 \Sh(Y_{\mathrm{proet}},A_{E,n})
 \longrightarrow\Sh(Y_{\mathrm{proet}},\mathcal O_E)
\]
is restriction of scalars along
\(\widehat{\mathcal O}_{E,Y}\twoheadrightarrow(A_{E,n})_Y\).  The reduction morphisms
\(A_{E,n+1}\twoheadrightarrow A_{E,n}\) give a tower in the
single derived category \(D(Y_{\mathrm{proet}},\mathcal O_E)\), represented
by a strict tower before applying \(R\varprojlim\).

Define \(D^b_c(Y,\mathcal O_E)\) to be the full
triangulated subcategory of derived-complete \(K\) such that
\[
 K_n:=K\otimes^L_{\mathcal O_E}A_{E,n}
 \in D^b_c(Y,A_{E,n})
\]
for every \(n\), and one interval \([r,s]\subset\mathbf Z\), independent of
\(n\), contains the cohomological support of every \(K_n\).
\uses{def:weil-small-proetale-site,def:weil-finite-torsion-constructible}
\end{definition}

\BlueprintLocalDeclarationRegistry{Definition}{def:weil-derived-complete-integral}{def:weil-small-proetale-site,def:weil-finite-torsion-constructible}{Lecture 41, Definitions 4.6 and 5.1 and Theorem 4.16.}

\begin{definition}[Rationalization and Karoubi completion]
\label{def:weil-rational-karoubi}
If \(\mathcal C\) is an \(\mathcal O_E\)-linear triangulated category, then
\(\mathcal C[1/\varpi_E]\) has the same objects and
\[
 \Hom_{\mathcal C[1/\varpi_E]}(K,L)
 =\Hom_{\mathcal C}(K,L)\otimes_{\mathcal O_E}E.
\]
The Karoubi completion \(\Kar(\mathcal C)\) has objects \((K,e)\), where
\(e^2=e\in\operatorname{End}_{\mathcal C}(K)\), and morphisms
\(u:(K,e)\to(L,e')\) satisfying \(u=e'ue\).  Set
\[
 D^b_c(Y,E)=
 \Kar\bigl(D^b_c(Y,\mathcal O_E)[1/\varpi_E]\bigr).
\]
For \(E=\Ql\), this uses \(\mathcal O_E=\Zl\) and
\(\varpi_E=\ell\).

If \(E\subset L\) is a finite extension, scalar extension is the exact
triangulated functor
\[
 -\otimes_E L:D^b_c(Y,E)\longrightarrow D^b_c(Y,L)
\]
obtained on integral representatives from completed derived tensor product
with \(\mathcal O_L\), levelwise modulo powers of the two uniformizers, and
then from rationalization and Karoubi completion.  Since
\(\mathcal O_L\) is finite free over \(\mathcal O_E\), this construction is
independent of an integral representative and is compatible with composition
of finite extensions.
\uses{def:weil-derived-complete-integral}
\end{definition}

\BlueprintLocalDeclarationRegistry{Definition}{def:weil-rational-karoubi}{def:weil-derived-complete-integral}{Lecture 41, Definitions 6.1, 6.4, and 6.6 and Theorems 6.3 and 6.5.}

\begin{definition}[Finite coefficient index and transition diagram]
\label{def:qbar-finext-diagram}
Let \(\mathcal E_\ell\) be the poset category whose objects are the finite
extensions \(\Ql\subset E\subset\Qlb\), with one arrow \(E\to L\) exactly
when \(E\subset L\).  The compositum of finitely many objects is a common
upper bound, so \(\mathcal E_\ell\) is filtered.  Put
\[
 \mathcal C_E(Y)=D^b_c(Y,E),
 \qquad
 T_{L/E}(K)=K\otimes_E L.
\]
The transition datum includes exact natural isomorphisms
\[
 \eta_E:T_{E/E}\xrightarrow{\sim}\operatorname{id},
 \qquad
 a_{M,L,E}:T_{M/L}T_{L/E}\xrightarrow{\sim}T_{M/E}
\]
for \(E\subset L\subset M\).  The unit equations, the tower pentagon, and
the shift comparisons are part of this datum.
\uses{conv:weil-size,def:weil-rational-karoubi}
\end{definition}

\BlueprintLocalDeclarationRegistry{Definition}{def:qbar-finext-diagram}{conv:weil-size,def:weil-rational-karoubi}{Lecture 41, Definitions 6.6 and 7.11 and Theorem 7.12; Lecture 42, Definition 1.3 and Theorem 1.4.}

\begin{definition}[Common-stage pair category and the symbol \(2\)-colimit]
\label{def:filtered-two-colimit}
An object is a pair \((E,K)\), where \(E\in\mathcal E_\ell\) and
\(K\in\mathcal C_E(Y)\).  A morphism representative
\[
 (M,u):(E,K)\longrightarrow(F,L)
\]
consists of a finite coefficient field \(M\) containing \(E\) and \(F\)
and a morphism
\[
 u:T_{M/E}K\longrightarrow T_{M/F}L
 \quad\text{in }\mathcal C_M(Y).
\]
If \(M\subset N\), its transport \(u_N\) is
\[
 T_{N/E}K\xrightarrow[\sim]{a^{-1}_{N,M,E}}
 T_{N/M}T_{M/E}K\xrightarrow{T_{N/M}(u)}
 T_{N/M}T_{M/F}L\xrightarrow[\sim]{a_{N,M,F}}T_{N/F}L.
\]
Representatives \((M,u)\) and \((N,v)\) are equivalent when a finite
coefficient field \(P\supset M,N\) satisfies \(u_P=v_P\).  Thus
\[
 \Hom((E,K),(F,L))=
 \varinjlim_{M\supset E,F}
 \Hom_{\mathcal C_M(Y)}(T_{M/E}K,T_{M/F}L),
\]
where the filtered colimit is taken in abelian groups and the displayed
equivalence relation is its representative description.

The identity is represented by the identity at stage \(E\), with the unit
isomorphism inserted.  Given \([M,u]:(E,K)\to(F,L)\) and
\([N,v]:(F,L)\to(G,Q)\), choose \(P\supset M,N\) and define
\[
 [N,v]\circ[M,u]=[P,v_P\circ u_P].
\]
Addition is defined by transporting two representatives to one common
stage and adding them there.  The zero object is \((\Ql,0)\); a biproduct of
\((E,K)\) and \((F,L)\) is represented over a common \(M\) by
\((M,T_{M/E}K\oplus T_{M/F}L)\).

The resulting additive category is denoted
\[
 \underset{E\in\mathcal E_\ell}{\operatorname{2-colim}^{\mathrm{pair}}}
 \mathcal C_E(Y).
\]
In this lecture the symbol \(\operatorname{2-colim}\) denotes this explicit
ordinary common-stage category, characterized by its pseudocone universal
property.  It is not an infinity-categorical colimit and it is not a union
of the object sets of the finite-stage categories.
\uses{def:qbar-finext-diagram}
\end{definition}

\BlueprintLocalDeclarationRegistry{Definition}{def:filtered-two-colimit}{def:qbar-finext-diagram}{Lecture 42, Definitions 2.1--2.11, Theorem 2.15, and Proposition 2.18.}

\begin{definition}[Common-stage shifts and distinguished triangles]
\label{def:qbar-pair-triangulation}
Define \((E,K)[1]=(E,K[1])\), with the transition shift isomorphisms used
on morphisms.  A triangle in the pair category is distinguished exactly
when it is isomorphic to the image of a distinguished triangle in
\(\mathcal C_M(Y)\) for one finite common stage \(M\).  For a representative
\((M,u)\), a cone triangle means the image of a chosen cone triangle for
\(u\) in \(\mathcal C_M(Y)\); no functorial choice of cones in a bare
triangulated category is asserted.  For \(\alpha\in\Qlb\), scalar
multiplication on a morphism is defined after enlarging to a common stage
containing \(\alpha\).
\uses{def:filtered-two-colimit}
\end{definition}

\BlueprintLocalDeclarationRegistry{Definition}{def:qbar-pair-triangulation}{def:filtered-two-colimit}{Lecture 42, Definitions 3.1, 3.3, 3.7, and 4.1 and Theorems 3.12 and 4.5.}

\begin{definition}[Algebraically closed coefficient category]
\label{def:qbar-coefficient-category}
For an additive category \(\mathcal C\), the Karoubi envelope has objects
\((A,e)\) with \(e^2=e\), morphisms
\[
 \Hom_{\Kar(\mathcal C)}((A,e),(B,d))
 =\{f:A\to B:f=dfe\},
\]
and identity \(e\) on \((A,e)\).  Shifts are \((A,e)[1]=(A[1],e[1])\),
and distinguished triangles are the triangles isomorphic to direct summands
of distinguished triangles in \(\mathcal C\).  Define
\[
 D^b_c(Y,\Qlb)=
 \Kar\left(
 \underset{E\in\mathcal E_\ell}{\operatorname{2-colim}^{\mathrm{pair}}}
 D^b_c(Y,E)\right).
\]
\uses{def:filtered-two-colimit,def:qbar-pair-triangulation}
\end{definition}

\BlueprintLocalDeclarationRegistry{Definition}{def:qbar-coefficient-category}{def:filtered-two-colimit,def:qbar-pair-triangulation}{Lecture 42, Definitions 5.1 and 5.3 and Theorem 5.4.}

\begin{theorem}[Correctness and universal property of the \(2\)-colimit]
\label{thm:qbar-two-colimit-model}
\uses{def:qbar-finext-diagram,def:filtered-two-colimit,
  def:qbar-pair-triangulation,def:qbar-coefficient-category}
The representative relation, addition, identities, and composition are
well defined.  The pair category is a \(\Qlb\)-linear ordinary triangulated
category with the stated shifts and distinguished triangles.  Its Karoubi
envelope is idempotent-complete and triangulated.

For every ordinary triangulated category \(\mathcal T\), exact functors from
the pair category to \(\mathcal T\) are equivalent to families of exact
functors \(F_E:\mathcal C_E(Y)\to\mathcal T\) equipped with unit-, tower-,
and shift-compatible isomorphisms
\(F_LT_{L/E}\xrightarrow{\sim}F_E\).  When \(\mathcal T\) is
idempotent-complete, the same universal property holds after Karoubi
completion.  Every object, morphism, finite diagram, and distinguished
triangle in \(D^b_c(Y,\Qlb)\) is isomorphic to one defined over a single
finite coefficient extension after enlarging that extension.
\sorryleaf{Lecture 42, Theorems 3.12, 4.5, 5.4, 6.4, and 6.5.}
\end{theorem}

\sketchslot{}

\begin{theorem}[Descent of compatible exact functors]
\label{thm:qbar-functor-descent}
\uses{thm:qbar-two-colimit-model}
Suppose exact functors \(\Phi_E:\mathcal C_E(X)\to\mathcal D_E(Y)\) are
equipped, for every \(E\subset L\), with exact natural isomorphisms
\[
 T^Y_{L/E}\Phi_E\xrightarrow{\sim}\Phi_LT^X_{L/E}
\]
satisfying the unit, tower, and shift equations.  Then they induce an exact
functor on pair categories and, after applying the idempotents, an exact
functor on their Karoubi envelopes.  Compatible natural transformations,
adjunction units and counits, base-change maps, projection maps, purity maps,
and duality maps descend by the same common-stage construction.  For
\(Rf^!\), the coefficient comparison is the mate of the comparison for
\(Rf_!\), and its tower coherence is part of the hypotheses.
\sorryleaf{Lecture 42, Definition 8.1 and Theorem 8.2.}
\end{theorem}

\sketchslot{}

\begin{longtable}{@{}
p{0.25\textwidth}
p{0.24\textwidth}
p{0.43\textwidth}@{}}
\caption{Components of the \(\Qlb\)-coefficient construction.}
\label{tab:qbar-two-colimit-sources}\\
\toprule
Component & Local declaration & Exact mathematical source\\
\midrule
\endfirsthead
\toprule
Component & Local declaration & Exact mathematical source\\
\midrule
\endhead
Filtered index and scalar diagram &
Definition~\ref{def:qbar-finext-diagram} &
Lecture 41, Definitions 6.6 and 7.11 and Theorem 7.12; Lecture 42,
Definition 1.3 and Theorem 1.4.\\
Objects, Hom representatives, equality, and composition &
Definition~\ref{def:filtered-two-colimit} &
Lecture 42, Definitions 2.1--2.11, Theorem 2.15, and Proposition 2.18.\\
Shifts, cones, triangles, and scalars &
Definition~\ref{def:qbar-pair-triangulation} &
Lecture 42, Definitions 3.1, 3.3, 3.7, and 4.1 and Theorems 3.12 and 4.5.\\
Karoubi triples &
Definition~\ref{def:qbar-coefficient-category} &
Lecture 42, Definitions 5.1 and 5.3 and Theorem 5.4.\\
Pseudocone universal property &
Theorem~\ref{thm:qbar-two-colimit-model} &
Lecture 42, Theorems 6.4 and 6.5.\\
Compatible exact functor descent &
Theorem~\ref{thm:qbar-functor-descent} &
Lecture 42, Definition 8.1 and Theorem 8.2.\\
\bottomrule
\end{longtable}

\begin{definition}[Finite Tor amplitude]
\label{def:finite-tor-amplitude}
For \(\Lambda=A_{E,n}\), an object
\(K\in D^b_c(Y,\Lambda)\) has \emph{finite Tor amplitude} if there are
integers \(a\leq b\) such that, for every sheaf \(M\) of
\(\Lambda\)-modules placed in degree zero,
\[
 \mathcal H^r(K\otimes^L_\Lambda M)=0
 \qquad\text{for }r\notin[a,b].
\]
For \(\Lambda=\mathcal O_E\), define finite Tor amplitude to mean that every
finite reduction has finite Tor amplitude in this sense, with one interval
independent of \(n\).  Write
\(D^b_{ctf}(Y,\Lambda)\) for the resulting full triangulated subcategory.
For field coefficients, define instead
\[
 D^b_{ctf}(Y,E)=D^b_c(Y,E),
 \qquad
 D^b_{ctf}(Y,\Qlb)=D^b_c(Y,\Qlb).
\]
The last equality is interpreted in the Karoubi-completed common-stage
category of Definition~\ref{def:qbar-coefficient-category}.
For \(\mathcal O_E\) and \(A_{E,n}\), finite Tor amplitude is an additional
condition and is never inferred merely from bounded constructibility.
\uses{def:weil-finite-torsion-constructible,def:weil-derived-complete-integral,
  def:weil-rational-karoubi,def:qbar-coefficient-category}
\end{definition}

\BlueprintLocalDeclarationRegistry{Definition}{def:finite-tor-amplitude}{def:weil-finite-torsion-constructible,def:weil-derived-complete-integral,def:weil-rational-karoubi,def:qbar-coefficient-category}{Lecture 40, Definition 9.6; Lecture 41, Proposition 3.8 and Corollary 5.2.}

\begin{definition}[Constructible cohomological intervals]
\label{def:constructible-interval}
For integers \(a\leq b\) and \(A=A_{E,n}\), let
\[
 D_c^{[a,b]}(Y,A)=
 \{K\in D^b_c(Y,A):\mathcal H^r(K)=0\text{ for }r\notin[a,b]\}.
\]
\uses{def:weil-finite-torsion-constructible}
\end{definition}

\BlueprintLocalDeclarationRegistry{Definition}{def:constructible-interval}{def:weil-finite-torsion-constructible}{Lecture 40, Definition 7.4; exact interval notation is local.}

\begin{definition}[Coefficient Tate twists]
\label{def:coefficient-twists}
Let \(S\) be a scheme on which \(\ell\) is invertible.  The pro-etale sheaf
\(\mu_{\ell^m,S}\) sends \(U\to S\) to
\[
 \{u\in\Gamma(U,\mathcal O_U^\times):u^{\ell^m}=1\}.
\]
With transition map \(u\mapsto u^\ell\), define
\[
 \Zl(1)_S=\varprojlim_m\mu_{\ell^m,S},\qquad
 \mathcal O_E(1)_S=\Zl(1)_S\otimes_{\Zl}\mathcal O_E,
\]
and
\[
 A_{E,n}(1)_S=\mathcal O_E(1)_S\otimes_{\mathcal O_E}A_{E,n},
 \qquad E(1)_S=\mathcal O_E(1)_S\otimes_{\mathcal O_E}E.
\]
The \(\Qlb\)-twist is represented by \(E(1)_S\) at any finite coefficient
stage and is compatible with scalar extension.  For every displayed
coefficient ring \(\Lambda\), put
\[
 \Lambda(m)=\Lambda(1)^{\otimes m},\qquad
 \Lambda(-m)=\operatorname{Hom}_\Lambda(\Lambda(m),\Lambda)
 \quad(m\geq0),
\]
and \(K(m)=K\otimes^L_\Lambda\Lambda(m)\).  Each \(\Lambda(m)\) is an
invertible locally constant rank-one sheaf, so twisting preserves
bounded constructibility and finite Tor amplitude.
\uses{def:weil-small-proetale-site,def:weil-finite-torsion-constructible,
  def:weil-derived-complete-integral,def:weil-rational-karoubi,
  def:qbar-coefficient-category,def:finite-tor-amplitude}
\end{definition}

\BlueprintLocalDeclarationRegistry{Definition}{def:coefficient-twists}{def:weil-small-proetale-site,def:weil-finite-torsion-constructible,def:weil-derived-complete-integral,def:weil-rational-karoubi,def:qbar-coefficient-category,def:finite-tor-amplitude}{Lecture 35, Definitions 2.1--2.4 and Proposition 2.5.}

\section{The six operations in triangulated form}

\begin{definition}[Derived pullback, tensor product, and internal Hom]
\label{def:derived-pullback-tensor-hom}
For a morphism \(f:X\to Y\) and a coefficient ring \(A=A_{E,n}\), define
\[
 Lf^*K=f^{-1}K\otimes^L_{f^{-1}A}A,
 \qquad
 -\otimes_A^L-,
 \qquad
 R\mathcal Hom_A(-,-)
\]
using K-flat resolutions for the derived tensor products and a K-injective
resolution of the second variable for derived internal Hom.  These
constructions give triangulated functors on the ambient derived categories.
For every \(n\), let \(f_n\) denote the morphism \(f\) equipped with
\(A_{E,n}\)-coefficients, and write
\[
 Lf_n^*=L(f_n)^*.
\]
For \(\mathcal O_E\), \(E\), and \(\Qlb\), they are obtained by compatible
finite reductions, derived inverse limit, rationalization, finite
coefficient extension, and Karoubi completion.
\uses{def:weil-small-proetale-site,def:weil-derived-complete-integral,
  def:weil-rational-karoubi,def:qbar-coefficient-category,
  thm:qbar-functor-descent}
\end{definition}

\BlueprintLocalDeclarationRegistry{Definition}{def:derived-pullback-tensor-hom}{def:weil-small-proetale-site,def:weil-derived-complete-integral,def:weil-rational-karoubi,def:qbar-coefficient-category,thm:qbar-functor-descent}{Lecture 40, Definitions 9.4 and 10.1 and Theorem 10.4; Lecture 41, Definitions 5.3 and 6.13.}

\begin{definition}[Derived direct image]
\label{def:derived-direct-image}
For \(f:X\to Y\), the direct-image functor on sheaves is left exact.  At a
finite level, its total right-derived functor is
\[
 Rf_{*,n}=R(f_n)_*:D(X_{\mathrm{proet}},A_{E,n})
 \longrightarrow D(Y_{\mathrm{proet}},A_{E,n}),
\]
computed by applying \(f_*\) to a K-injective resolution.  For a
derived-complete integral complex \(K\), define
\[
 Rf_*K=
 R\varprojlim_n\operatorname{Res}_n Rf_{*,n}
 \bigl(K\otimes^L_{\mathcal O_E}A_{E,n}\bigr),
\]
where the tower maps are induced by
\(A_{E,n+1}\twoheadrightarrow A_{E,n}\).  The finite-reduction comparison is
\[
 (Rf_*K)\otimes^L_{\mathcal O_E}A_{E,n}
 \xrightarrow{\sim}
 Rf_{*,n}\bigl(K\otimes^L_{\mathcal O_E}A_{E,n}\bigr).
\]
The rational and \(\Qlb\)-valued ambient functors are induced at a finite
coefficient stage and then passed through rationalization and Karoubi
completion.
\uses{def:weil-small-proetale-site,def:weil-derived-complete-integral,
  def:weil-rational-karoubi,def:qbar-coefficient-category,
  thm:qbar-functor-descent}
\end{definition}

\BlueprintLocalDeclarationRegistry{Definition}{def:derived-direct-image}{def:weil-small-proetale-site,def:weil-derived-complete-integral,def:weil-rational-karoubi,def:qbar-coefficient-category,thm:qbar-functor-descent}{Lecture 30, Definition 2.5; Lecture 43, Definition 3.1.}

\begin{definition}[Extension by zero]
\label{def:extension-by-zero}
Let \(A=A_{E,n}\).  For an open immersion
\(j:U\hookrightarrow\overline X\), extension by zero is the exact functor
\[
 j_!: \Sh(U_{\mathrm{proet}},A)
 \longrightarrow \Sh(\overline X_{\mathrm{proet}},A)
\]
left adjoint to \(j^*\), characterized on geometric stalks by
\[
 (j_!\mathcal F)_{\bar x}=
 \begin{cases}
  \mathcal F_{\bar x},&\bar x\text{ lies in }U,\\
  0,&\bar x\text{ lies in }\overline X\setminus U.
 \end{cases}
\]
Exactness gives, by degreewise application, an exact triangulated functor
\[
 j_!:D(U_{\mathrm{proet}},A)
 \longrightarrow D(\overline X_{\mathrm{proet}},A).
\]
\uses{def:weil-small-proetale-site,def:weil-finite-torsion-constructible}
\end{definition}

\BlueprintLocalDeclarationRegistry{Definition}{def:extension-by-zero}{def:weil-small-proetale-site,def:weil-finite-torsion-constructible}{Lecture 40, Definition 5.11; Lecture 43, Definitions 2.4 and 2.13.}

\begin{definition}[Compactifications and refinements]
\label{def:compactification-category}
Let \(f:X\to Y\) be separated and of finite type.  A
\emph{compactification} \(c=(j,\bar f)\) of \(f\) is a factorization
\[
 X\xrightarrow{j}\overline X\xrightarrow{\bar f}Y
\]
with \(j\) an open immersion, \(\bar f\) proper, and
\(\bar f\circ j=f\).  A \emph{refinement}
\(u:c_1\to c_2\) is a proper \(Y\)-morphism
\[
 u:\overline X_1\longrightarrow\overline X_2
\]
such that \(u\circ j_1=j_2\).  Compactifications and refinements form a
category \(\operatorname{Comp}(f)\).
\uses{def:weil-geometric-hypotheses}
\end{definition}

\BlueprintLocalDeclarationRegistry{Definition}{def:compactification-category}{def:weil-geometric-hypotheses}{Lecture 43, Definitions 4.1 and 4.3.}

\begin{proposition}[Existence and common refinement]
\label{prop:compactification-refinement}
\uses{def:compactification-category}
The category \(\operatorname{Comp}(f)\) is nonempty.  Given finitely many
compactifications of \(f\), there is a compactification refining all of
them.  Given two parallel refinements, there is a further refinement after
which they agree.
\end{proposition}

\sketchslot{}


\begin{definition}[Compact-support candidate attached to a compactification]
\label{def:compact-support-candidate}
Let \(A=A_{E,n}\), let \(c=(j,\bar f)\in\operatorname{Comp}(f)\), and set
\[
 Rf_{!,n}^{c}=R\bar f_{*,n}\,j_!:
 D(X_{\mathrm{proet}},A)
 \longrightarrow D(Y_{\mathrm{proet}},A).
\]
Here \(R\bar f_{*,n}\) is the K-injective derived direct image and \(j_!\)
is the exact derived functor of extension by zero.
\uses{def:derived-direct-image,def:extension-by-zero,
  def:compactification-category}
\end{definition}

\BlueprintLocalDeclarationRegistry{Definition}{def:compact-support-candidate}{def:derived-direct-image,def:extension-by-zero,def:compactification-category}{Lecture 43, Definition 4.11.}

\begin{proposition}[Compactification coherence]
\label{prop:compactification-coherence}
\uses{prop:compactification-refinement,def:compact-support-candidate}
For every refinement \(u:c_1\to c_2\), the canonical transformation
\[
 R\bar f_{1*,n}j_{1!}
 \xrightarrow{\sim}
 R\bar f_{2*,n}j_{2!}
\]
is the composite of the proper-composition isomorphism with the isomorphism
\[
 Ru_*j_{1!}\xrightarrow{\sim}j_{2!}
\]
on the ambient derived category.  These transformations are identities for
identity refinements and are transitive under composition.  For any two
compactifications \(c_1,c_2\), comparison through a common refinement gives
a natural isomorphism
\[
 Rf_{!,n}^{c_1}\xrightarrow{\sim}Rf_{!,n}^{c_2}.
\]
It is independent of the common refinement and satisfies the cocycle
condition for triples of compactifications.
\end{proposition}

\sketchslot{}


\begin{definition}[Compactification-independent finite-level direct image]
\label{def:compact-support-direct-image}
Choose \(c\in\operatorname{Comp}(f)\).  Define
\[
 Rf_{!,n}=Rf_{!,n}^{c}:
 D(X_{\mathrm{proet}},A_{E,n})
 \longrightarrow D(Y_{\mathrm{proet}},A_{E,n}),
\]
equipped with the canonical comparison isomorphisms of
Proposition~\ref{prop:compactification-coherence}.  The cocycle condition
makes this an exact triangulated functor independent of \(c\) up to the
specified unique coherent isomorphism.
\uses{def:compact-support-candidate,prop:compactification-coherence}
\end{definition}

\BlueprintLocalDeclarationRegistry{Definition}{def:compact-support-direct-image}{def:compact-support-candidate,prop:compactification-coherence}{Lecture 43, Definition 5.7 and Theorem 5.8.}

\begin{proposition}[Finite-level bounded constructibility for proper support]
\label{prop:finite-compact-support-constructibility}
\uses{def:compact-support-direct-image,def:constructible-interval}
For every separated finite-type \(f:X\to Y\), the ambient functor
\(Rf_{!,n}\) restricts to an exact triangulated functor
\[
 Rf_{!,n}:D^b_c(X,A_{E,n})
 \longrightarrow D^b_c(Y,A_{E,n}).
\]
For each cohomological interval \([a,b]\) and each upper bound on the
relative dimension of \(f\), there is an interval \([a',b']\), independent
of the compactification, such that
\[
 K\in D_c^{[a,b]}(X,A_{E,n})
 \quad\Longrightarrow\quad
 Rf_{!,n}K\in D_c^{[a',b']}(Y,A_{E,n}).
\]
The restriction and its amplitude bound commute with finite coefficient
reduction.
\end{proposition}

\sketchslot{}


\begin{definition}[Integral, rational, and algebraically closed
compact-support functors]
\label{def:all-coefficient-compact-support}
For \(K\in D^b_c(X,\mathcal O_E)\), put
\(K_n=K\otimes^L_{\mathcal O_E}A_{E,n}\) and define
\[
 Rf_!K=
 R\varprojlim_n\operatorname{Res}_n Rf_{!,n}K_n.
\]
The \(E\)-coefficient functor is induced from this functor by
rationalization and Karoubi completion.  For a finite extension \(E\subset
L\), its scalar-extension comparison is represented by the finite-level
comparisons over \(\mathcal O_E\subset\mathcal O_L\).  Passing to the
common finite coefficient stage and then taking the Karoubi completion
defines the \(\Qlb\)-coefficient functor.
\uses{def:weil-derived-complete-integral,def:weil-rational-karoubi,
  def:qbar-coefficient-category,def:compact-support-direct-image,
  thm:qbar-functor-descent}
\end{definition}

\BlueprintLocalDeclarationRegistry{Definition}{def:all-coefficient-compact-support}{def:weil-derived-complete-integral,def:weil-rational-karoubi,def:qbar-coefficient-category,def:compact-support-direct-image,thm:qbar-functor-descent}{Exact local packaging; constructible realization is Lecture 43, Theorem 9.17.}

\begin{proposition}[Coefficient passage and constructible landing for
proper support]
\label{prop:compact-support-coefficient-passage}
\uses{prop:finite-compact-support-constructibility,
  def:all-coefficient-compact-support,thm:qbar-functor-descent}
For \(\Lambda\in\{\mathcal O_E,E,\Qlb\}\), the preceding construction is
an exact triangulated functor
\[
 Rf_!:D^b_c(X,\Lambda)\longrightarrow D^b_c(Y,\Lambda).
\]
For integral coefficients there are natural isomorphisms
\[
 (Rf_!K)\otimes^L_{\mathcal O_E}A_{E,n}
 \xrightarrow{\sim}Rf_{!,n}K_n.
\]
These isomorphisms, the compactification comparisons, and composition are
compatible with reduction, rationalization, finite scalar extension, and
the \(\Qlb\)-valued common-stage construction.
\end{proposition}

\sketchslot{}


\begin{definition}[Full cohomological dimension]
\label{def:full-cohomological-dimension}
Let \(A=A_{E,n}\), and let \(p:\overline X\to Y\) be proper.  Write
\(\operatorname{FullCD}(p,A;c)\) if an integer
\(c\geq0\) satisfies
\[
 R^qp_*\mathcal F=0\qquad(q>c)
\]
for every sheaf \(\mathcal F\) of \(A\)-modules, with the same bound after
restriction over every object of the small pro-etale site of \(Y\).
\uses{def:weil-small-proetale-site,def:derived-direct-image}
\end{definition}

\BlueprintLocalDeclarationRegistry{Definition}{def:full-cohomological-dimension}{def:weil-small-proetale-site,def:derived-direct-image}{Lecture 44, Definition 3.7.}

\begin{proposition}[Ambient exceptional right adjoint]
\label{prop:ambient-exceptional-adjoint}
\uses{def:full-cohomological-dimension,def:compact-support-direct-image,
  prop:compactification-coherence}
Let \(A=A_{E,n}\).  Let \(f:X\to Y\) be separated and of finite type, and
choose a compactification
\(X\xrightarrow{j}\overline X\xrightarrow{p}Y\), and suppose that
\(\operatorname{FullCD}(p,A;c)\) holds.  The exact functor
\[
 Rf_{!,n}:D(X_{\mathrm{proet}},A)
 \longrightarrow D(Y_{\mathrm{proet}},A)
\]
preserves arbitrary coproducts and has a right adjoint
\[
 Rf_n^{\times}:D(Y_{\mathrm{proet}},A)
 \longrightarrow D(X_{\mathrm{proet}},A).
\]
The adjunction is independent of the compactification under the coherent
comparisons of Proposition~\ref{prop:compactification-coherence}.
\end{proposition}

\sketchslot{}


\begin{definition}[Exceptional-adjoint admissibility]
\label{def:exceptional-adjoint-admissible}
Let \(A=A_{E,n}\), let \(f:X\to Y\) be separated and of finite type, and
choose a compactification
\(X\xrightarrow{j}\overline X\xrightarrow{p}Y\).  Say that \(f\) is
\emph{exceptional-adjoint admissible at \(A\)} if:
\begin{enumerate}
\item[\textup{(a)}] \(\operatorname{FullCD}(p,A;c)\) holds for some \(c\), and
\(Rf_n^{\times}\) is the ambient right adjoint furnished by
Proposition~\ref{prop:ambient-exceptional-adjoint};
\item[\textup{(b)}] there is an integer \(c_f\) such
that
\[
 L\in D_c^{[a,b]}(Y,A)\Longrightarrow
 Rf_n^{\times}L\in D_c^{[a-c_f,b]}(X,A)
\]
for all \(a\leq b\);
\item[\textup{(c)}] one has
\[
 L\in D^b_{ctf}(Y,A)\Longrightarrow
 Rf_n^{\times}L\in D^b_{ctf}(X,A).
\]
\end{enumerate}
At the \(\mathcal O_E\), \(E\), or \(\Qlb\) stage, admissibility means that
these hypotheses hold at the finite stages with uniform bounds and that the
finite-reduction and scalar-extension comparisons for the representing
objects are isomorphisms.  Admissibility is always an explicit hypothesis;
it is not inferred from separated finite type alone.
\uses{def:full-cohomological-dimension,prop:ambient-exceptional-adjoint,
  def:constructible-interval,def:finite-tor-amplitude,
  def:weil-derived-complete-integral,def:weil-rational-karoubi,
  def:qbar-coefficient-category}
\end{definition}

\BlueprintLocalDeclarationRegistry{Definition}{def:exceptional-adjoint-admissible}{def:full-cohomological-dimension,prop:ambient-exceptional-adjoint,def:constructible-interval,def:finite-tor-amplitude,def:weil-derived-complete-integral,def:weil-rational-karoubi,def:qbar-coefficient-category}{Exact local admissibility package; the resulting formalism is Lecture 44, Theorem 6.20.}

\begin{definition}[Exceptional inverse image notation]
\label{def:exceptional-inverse-image}
At a finite stage at which \(f\) is exceptional-adjoint admissible, set
\[
 Rf_n^!=Rf_n^{\times}|_{D^b_c(Y,A_{E,n})}:
 D^b_c(Y,A_{E,n})\longrightarrow D^b_c(X,A_{E,n}).
\]
For \(\Lambda\in\{\mathcal O_E,E,\Qlb\}\), at a stage at which \(f\) is
admissible, the compatible finite-reduction right adjoints, followed
respectively by derived inverse limit, rationalization, or the common-stage
Karoubi construction, define
\[
 Rf^!:D^b_c(Y,\Lambda)\longrightarrow D^b_c(X,\Lambda).
\]
In every case this is right adjoint to the previously defined restriction of
\(Rf_!\), and it is characterized by natural isomorphisms
\[
 \Hom_Y(Rf_!K,L)\cong\Hom_X(K,Rf^!L).
\]
Its restriction to finite-Tor objects is part of admissibility.  Smooth
purity and Poincare duality for smooth proper structure morphisms are stated
independently below.
\uses{def:exceptional-adjoint-admissible,prop:ambient-exceptional-adjoint,
  def:all-coefficient-compact-support,thm:qbar-functor-descent}
\end{definition}

\BlueprintLocalDeclarationRegistry{Definition}{def:exceptional-inverse-image}{def:exceptional-adjoint-admissible,prop:ambient-exceptional-adjoint,def:all-coefficient-compact-support,thm:qbar-functor-descent}{Lecture 44, Definition 6.19 and Theorem 6.20.}

\begin{definition}[Exchange and projection transformations]
\label{def:exchange-projection}
Consider a Cartesian square
\[
\begin{array}{ccc}
X' & \xrightarrow{g'} & X\\
{\scriptstyle f'}\downarrow && \downarrow{\scriptstyle f}\\
Y' & \xrightarrow{g} & Y.
\end{array}
\]
The direct-image exchange transformation
\[
 \operatorname{Ex}_*:\;Lg^*Rf_*\longrightarrow Rf'_*Lg'^*
\]
is adjoint to the composite
\[
 Lf'^*Lg^*Rf_*\xrightarrow{\sim}Lg'^*Lf^*Rf_*
 \longrightarrow Lg'^*,
\]
whose second arrow is the counit of \(Lf^*\dashv Rf_*\).  If
\(f=\bar f\circ j\) is a compactification, the compact-support exchange map
\[
 \operatorname{Ex}_!:\;Lg^*Rf_!\longrightarrow Rf'_!Lg'^*
\]
is represented by \(\operatorname{Ex}_*\) for \(\bar f\), followed by the
canonical identification
\(L\bar g'^*j_!\cong j'_!Lg'^*\); refinement coherence makes it independent
of the compactification.

For \(K\in D^b_c(X,\Lambda)\) and
\(L\in D^b_{ctf}(Y,\Lambda)\), the projection transformation is
\[
 \operatorname{PF}_{f,K,L}:Rf_!K\otimes^L L
 \longrightarrow Rf_!(K\otimes^L Lf^*L).
\]
For a compactification it is the direct-image projection transformation for
\(\bar f\), followed by
\(j_!K\otimes^LL\bar f^*L\cong
j_!(K\otimes^LLf^*L)\).  These are natural transformations with the
displayed sources and targets.
\uses{def:derived-pullback-tensor-hom,def:derived-direct-image,
  def:all-coefficient-compact-support,def:compactification-category,
  prop:compactification-coherence}
\end{definition}

\BlueprintLocalDeclarationRegistry{Definition}{def:exchange-projection}{def:derived-pullback-tensor-hom,def:derived-direct-image,def:all-coefficient-compact-support,def:compactification-category,prop:compactification-coherence}{Lecture 31, Definitions 1.13 and 11.12; Lecture 33, Definitions 2.1 and 5.4.}

\begin{proposition}[Geometric fibers and algebraically closed extension]
\label{prop:proper-fibers}
\uses{def:derived-direct-image,def:exchange-projection}
Let \(A=A_{E,n}\).  Let \(f:X\to Y\) be proper between finite-type schemes
over an algebraically closed field of characteristic different from \(\ell\),
and let \(K\in D^b_c(X,A)\).
\begin{enumerate}
\item[\textup{(a)}] If \(\Omega\subset\Omega'\) is an extension of algebraically closed
fields, \(Z/\Omega\) is proper, and \(L\in D^b_c(Z,A)\), then base extension
induces an isomorphism
\[
 Lg^*R(a_Z)_*L\xrightarrow{\sim}
 R(a_{Z_{\Omega'}})_*Lg'^*L
\]
in \(D^b_c(\Spec\Omega',A)\).
\item[\textup{(b)}] For every geometric point \(\bar y\to Y\), with geometric fiber
\(X_{\bar y}=X\times_Y\bar y\), the exchange morphism is the fiber
isomorphism
\[
 L\bar y^*Rf_*K\xrightarrow{\sim}
 R(a_{X_{\bar y}})_*L\bar g'^*K.
\]
\end{enumerate}
\end{proposition}

\sketchslot{}


\begin{proposition}[Finite-torsion proper base change and constructibility]
\label{prop:finite-proper-base-change}
\uses{def:exchange-projection,prop:proper-fibers,
  def:constructible-interval}
Let \(A=A_{E,n}\), and let the square in
Definition~\ref{def:exchange-projection} be an arbitrary Cartesian
square of finite-type schemes on which \(\ell\) is invertible.  If \(f\) is
proper, then for every \(K\in D^b_c(X,A)\),
\[
 \operatorname{Ex}_*(K):Lg^*Rf_*K\xrightarrow{\sim}Rf'_*Lg'^*K.
\]
Moreover, \(Rf_*K\) is bounded constructible.  More explicitly, each
\(R^jf_*\mathcal H^r(K)\) is constructible, and if the geometric fibers of
\(f\) have dimension at most \(e\) and \(\mathcal H^r(K)=0\) outside
\([u,v]\), then \(\mathcal H^m(Rf_*K)=0\) outside \([u,v+2e]\).
The exchange isomorphism is compatible with composition of Cartesian
squares and specializes at every geometric point to
Proposition~\ref{prop:proper-fibers}.
\end{proposition}

\sketchslot{}


\begin{proposition}[Adic, rational, and algebraically closed passage]
\label{prop:coefficient-passage}
\uses{prop:finite-proper-base-change,def:derived-direct-image,
  def:weil-derived-complete-integral,def:weil-rational-karoubi,
  thm:qbar-functor-descent}
Let \(f:X\to Y\) be proper and let
\(K\in D^b_c(X,\mathcal O_E)\).  Put
\(K_n=K\otimes^L_{\mathcal O_E}A_{E,n}\).  Then
\[
 (Rf_*K)\otimes^L_{\mathcal O_E}A_{E,n}
 \xrightarrow{\sim}Rf_{*,n}K_n
\]
for every \(n\), and
\[
 Rf_*K\xrightarrow{\sim}
 R\varprojlim_n\operatorname{Res}_nRf_{*,n}K_n.
\]
The reduction of \(\operatorname{Ex}_*(K)\) is
\(\operatorname{Ex}_*(K_n)\).  At each coefficient level
\(\Lambda\in\{A_{E,n},\mathcal O_E,E,\Qlb\}\),
\[
 Rf_*:D^b_c(X,\Lambda)\longrightarrow D^b_c(Y,\Lambda)
\]
is a well-defined exact triangulated functor and \(\operatorname{Ex}_*\) is
an isomorphism for proper \(f\).
\end{proposition}

\sketchslot{}


\begin{proposition}[Coefficient and functor package]
\label{prop:functor-package}
\uses{def:derived-pullback-tensor-hom,def:derived-direct-image,
  def:all-coefficient-compact-support,def:exceptional-inverse-image,
  def:exchange-projection,prop:compact-support-coefficient-passage,
  prop:coefficient-passage,thm:qbar-functor-descent}
Let \(f:X\to Y\) be a separated morphism of finite type between schemes of
finite type over \(k\).  Let
\[
 \Lambda\in\{A_{E,n},\mathcal O_E,E,\Qlb\}.
\]
Then the ambient constructions above preserve bounded constructibility and
give exact triangulated functors
\[
 \begin{aligned}
 Lf^*&:D^b_c(Y,\Lambda)\longrightarrow D^b_c(X,\Lambda),\\
 Rf_*,Rf_!&:D^b_c(X,\Lambda)\longrightarrow D^b_c(Y,\Lambda),\\
 -\otimes^L- &:D^b_{ctf}(Y,\Lambda)\times D^b_c(Y,\Lambda)
              \longrightarrow D^b_c(Y,\Lambda),\\
 R\mathcal Hom(-,-)&:D^b_{ctf}(Y,\Lambda)^{\mathrm{op}}
              \times D^b_c(Y,\Lambda)
              \longrightarrow D^b_c(Y,\Lambda).
 \end{aligned}
\]
If \(f\) is exceptional-adjoint admissible at \(\Lambda\), there is in
addition an exact functor
\[
 Rf^!:D^b_c(Y,\Lambda)\longrightarrow D^b_c(X,\Lambda),
\]
and the adjunction \(Rf_!\dashv Rf^!\) restricts to
\[
 Rf_!:D^b_{ctf}(X,\Lambda)
 \rightleftarrows D^b_{ctf}(Y,\Lambda):Rf^!.
\]
They satisfy the following statements.
\begin{enumerate}
\item[\textup{(a)}] The scalar-extension functors between finite coefficient fields are
exact and preserve shifts and distinguished triangles.
The common-stage construction of \(\mathcal C_Y\) satisfies the triangulated
axioms, and its Karoubi completion is triangulated.  Every finite-stage exact
functor and natural isomorphism compatible with scalar extension descends to
\(D^b_c(Y,\Qlb)\).
\item[\textup{(b)}] The functor \(Rf_!\) and its canonical comparison maps are those of
Definition~\ref{def:compact-support-direct-image},
Proposition~\ref{prop:compactification-coherence}, and
Proposition~\ref{prop:compact-support-coefficient-passage}.
\item[\textup{(c)}] For composable separated finite-type morphisms,
\[
 R(g\circ f)_!\xrightarrow{\sim}Rg_!Rf_!,
\]
with identity compatibility and the associativity pentagon.  If \(f\) is
proper, \(Rf_!\xrightarrow{\sim}Rf_*\).
For a closed immersion \(i:Z\hookrightarrow X\) with open complement
\(j:U\hookrightarrow X\), there is a functorial distinguished triangle
\[
 R(fj)_!j^*K\longrightarrow Rf_!K\longrightarrow
 R(fi)_!i^*K\longrightarrow R(fj)_!j^*K[1].
\]
\item[\textup{(d)}] In a Cartesian square
\[
\begin{array}{ccc}
X' & \xrightarrow{g'} & X\\
{\scriptstyle f'}\downarrow && \downarrow{\scriptstyle f}\\
Y' & \xrightarrow{g} & Y,
\end{array}
\]
 the exchange map of Definition~\ref{def:exchange-projection}
\(Lg^*Rf_*\to Rf'_*Lg'^*\) is an isomorphism if \(f\) is proper or if
\(g\) is smooth.  For arbitrary \(g\),
\[
 Lg^*Rf_!\xrightarrow{\sim}Rf'_!Lg'^*.
\]
\item[\textup{(e)}] The projection transformation is an isomorphism:
\[
 Rf_!K\otimes^L L\xrightarrow{\sim}Rf_!(K\otimes^L Lf^*L).
\]
Here \(K\in D^b_c(X,\Lambda)\) and
\(L\in D^b_{ctf}(Y,\Lambda)\).
\item[\textup{(f)}] Assume \(f\) is exceptional-adjoint admissible.  If \(f\) is smooth
of pure relative dimension \(r\) and
\(L\in D^b_{ctf}(Y,\Lambda)\), then
\[
 Rf^!L\xrightarrow{\sim}Lf^*L(r)[2r].
\]
\item[\textup{(g)}] Assume the structure morphisms and \(f\) are exceptional-adjoint
admissible.  Put \(\omega_Y=Ra_Y^!\Lambda\) and
\(D_YK=R\mathcal Hom(K,\omega_Y)\).  On
\(D^b_{ctf}(Y,\Lambda)\), for the stated finite-type \(Y/k\) and
\(\ell\neq\operatorname{char}(k)\), \(D_Y\) is a
contravariant exact equivalence and \(K\to D_YD_YK\) is an isomorphism.  For
\(K\in D^b_{ctf}(X,\Lambda)\),
\[
 D_YRf_!K\xrightarrow{\sim}Rf_*D_XK.
\]
\end{enumerate}
For the coefficient comparisons, use the finite-stage symbols
\(Lf_n^*\), \(Rf_{*,n}\), \(Rf_{!,n}\), and \(Rf_n^!\) introduced with
the corresponding operations above.
For an exceptional-adjoint admissible structure morphism
\(a_Y:Y\to\Spec k\), put
\[
 \omega_{Y,n}=Ra_{Y,n}^!A_{E,n},\qquad
 D_{Y,n}(-)=R\mathcal Hom_{A_{E,n}}(-,\omega_{Y,n}).
\]
For
\(K\in D^b_{ctf}(X,\mathcal O_E)\),
\(M\in D^b_{ctf}(Y,\mathcal O_E)\), and
\(N\in D^b_c(Y,\mathcal O_E)\), write
\(K_n=K\otimes^L_{\mathcal O_E}A_{E,n}\), and similarly for \(M_n,N_n\).
There are natural isomorphisms
\[
\begin{aligned}
 (Lf^*N)_n&\xrightarrow{\sim}Lf_n^*N_n,&
 (Rf_*K)_n&\xrightarrow{\sim}Rf_{*,n}K_n,\\
 (Rf_!K)_n&\xrightarrow{\sim}Rf_{!,n}K_n,&
 (Rf^!M)_n&\xrightarrow{\sim}Rf_n^!M_n,\\
 (M\otimes^LN)_n&\xrightarrow{\sim}M_n\otimes^L_{A_{E,n}}N_n,&
 \bigl(R\mathcal Hom(M,N)\bigr)_n&\xrightarrow{\sim}
 R\mathcal Hom_{A_{E,n}}(M_n,N_n),\\
 (D_YM)_n&\xrightarrow{\sim}D_{Y,n}M_n.&
\end{aligned}
\]
The comparison involving \(Rf^!\) assumes that \(f\) is
exceptional-adjoint admissible; the one involving \(D_Y\) assumes the same
for \(a_Y\).  If \(f\) is also smooth of relative dimension \(r\), reduction
identifies the
map
\((Rf^!M)_n\to(Lf^*M(r)[2r])_n\) with the finite-level purity map
\[
 Rf_n^!M_n\xrightarrow{\sim}Lf_n^*M_n(r)[2r].
\]
Reduction also carries \(\operatorname{Ex}_*\), \(\operatorname{Ex}_!\),
\(\operatorname{PF}\), biduality, and
\(D_YRf_!K\xrightarrow{\sim}Rf_*D_XK\) to their finite-level
counterparts.  The displayed finite-Tor hypotheses are required precisely
for the tensor, internal-Hom, exceptional-pullback, purity, and duality
comparisons.  All these finite-stage functors and transformations commute
with extension between finite coefficient fields and induce the
stated \(\Qlb\)-valued operations.
\end{proposition}

\sketchslot{}


\section{Weil structures, cohomology, and Frobenius factors}

\begin{proposition}[Construction of constructible truncations]
\label{prop:constructible-truncations}
\uses{def:qbar-coefficient-category,thm:qbar-two-colimit-model}
There are truncation functors
\[
 \tau^{\leq m},\tau^{\geq m}:D^b_c(Y,\Qlb)
 \longrightarrow D^b_c(Y,\Qlb).
\]
They are represented at a common finite coefficient field by the usual
constructible truncations and commute with scalar extension and idempotent
summands.
For \(\Lambda=E\) or \(\Qlb\), these truncations, their natural maps, and the
distinguished truncation triangles satisfy the axioms of a bounded
t-structure.  Define its cohomology functors by
\[
 \mathcal H^i(K)=\bigl(\tau^{\geq i}\tau^{\leq i}K\bigr)[i].
\]
Pullback by a scheme automorphism commutes with the truncations and with
\(\mathcal H^i\).
\end{proposition}

\sketchslot{}


\begin{definition}[Constructible heart]
\label{def:constructible-heart}
For \(\Lambda=E\) or \(\Qlb\), define
\[
\begin{aligned}
 D^{\leq0}_c(Y,\Lambda)&=
 \{K:\tau^{\geq1}K=0\},\\
 D^{\geq0}_c(Y,\Lambda)&=
 \{K:\tau^{\leq-1}K=0\}.
\end{aligned}
\]
Their intersection is the abelian category \(\Sh_c(Y,\Lambda)\), called the
constructible heart.  Its objects are the constructible \(\Lambda\)-adic
sheaves.
\uses{prop:constructible-truncations,def:qbar-coefficient-category}
\end{definition}

\BlueprintLocalDeclarationRegistry{Definition}{def:constructible-heart}{prop:constructible-truncations,def:qbar-coefficient-category}{Lecture 40, Theorem 7.7; Lecture 42, Theorem 5.8.}

\begin{definition}[Weil sheaf and Weil complex]
\label{def:strict-weil-complex}
Let \(Y_0\) be of finite type over \(k_0\), with geometric fiber \(Y\).  A
constructible \(E\)-Weil sheaf is an object
\(\mathcal F\in\Sh_c(Y,E)\) together with an isomorphism
\[
 \theta:Lg_Y^*\mathcal F\xrightarrow{\sim}\mathcal F.
\]
Morphisms commute with \(\theta\).  The category
\(D^b_{c,\mathrm W}(Y,E)\) is obtained from bounded complexes of
constructible \(E\)-Weil sheaves by localizing at equivariant
quasi-isomorphisms.  It is an ordinary triangulated category.  Define
\[
 D^b_{c,\mathrm W}(Y,\Qlb)=
 \Kar\left(\underset{E\in\mathcal E_\ell}{\operatorname{2-colim}}
 D^b_{c,\mathrm W}(Y,E)\right)
\]
using the same common-stage description as for \(D^b_c(Y,\Qlb)\).
For either coefficient field there is an exact forgetful functor
\[
 \operatorname{For}:D^b_{c,\mathrm W}(Y,\Lambda)
 \longrightarrow D^b_c(Y,\Lambda).
\]
An object \(K\) carries, in addition to \(\operatorname{For}(K)\), the
derived Frobenius isomorphism
\[
 \theta_K:Lg_Y^*\operatorname{For}(K)
 \xrightarrow{\sim}\operatorname{For}(K),
\]
represented termwise at a finite coefficient stage.
\uses{def:weil-base-field,def:weil-frobenius-convention,
  def:constructible-heart,def:qbar-finext-diagram,
  def:filtered-two-colimit,def:qbar-coefficient-category,
  thm:qbar-two-colimit-model}
\end{definition}

\BlueprintLocalDeclarationRegistry{Definition}{def:strict-weil-complex}{def:weil-base-field,def:weil-frobenius-convention,def:constructible-heart,def:qbar-finext-diagram,def:filtered-two-colimit,def:qbar-coefficient-category,thm:qbar-two-colimit-model}{Lecture 42, Definitions 9.1, 9.3, 10.1, and 10.26 and Theorem 10.29.}

\begin{proposition}[Six operations on Weil complexes]
\label{prop:weil-operations}
\uses{prop:functor-package,prop:constructible-truncations,
  def:strict-weil-complex,thm:qbar-functor-descent}
Let \(f_0:X_0\to Y_0\) be separated and of finite type, and let
\(f:X\to Y\) be its geometric fiber.  The exact functors
\[
 Lf^*,\quad Rf_*,\quad Rf_!,\quad -\otimes^L-,\quad R\mathcal Hom
\]
on the domains in Proposition~\ref{prop:functor-package} lift to Weil
complexes.  If \(f\) is exceptional-adjoint admissible, the same holds for
\(Rf^!\).  Explicitly, the Weil isomorphism on \(Rf_*K\) is
\[
 Lg_Y^*Rf_*\operatorname{For}(K)
 \xrightarrow{\operatorname{Ex}_*}
 Rf_*Lg_X^*\operatorname{For}(K)
 \xrightarrow{Rf_*(\theta_K)}Rf_*\operatorname{For}(K),
\]
and the compact-support construction uses \(\operatorname{Ex}_!\).
The forgetful functor commutes with these lifted operations.  Composition,
base change, projection, purity, and duality isomorphisms are morphisms of
Weil complexes and induce Frobenius-equivariant maps on cohomology.
\end{proposition}

\sketchslot{}


\begin{definition}[Frobenius on coefficient twists]
\label{def:frobenius-on-twists}
The coefficient twists of Definition~\ref{def:coefficient-twists} on
\(\Spec k_0\), pulled back to a
geometric fiber, have their natural Weil isomorphisms.  Geometric Frobenius
acts on \(E(1)\) by \(q^{-1}\), on \(E(m)\) by \(q^{-m}\), and on
\(E(-m)\) by \(q^m\).  These assertions commute with finite reduction and
extension to \(\Qlb\); in particular, geometric Frobenius acts on
\(E(-d)\) by \(q^d\).
\uses{def:weil-frobenius-convention,def:coefficient-twists,
  def:strict-weil-complex}
\end{definition}

\BlueprintLocalDeclarationRegistry{Definition}{def:frobenius-on-twists}{def:weil-frobenius-convention,def:coefficient-twists,def:strict-weil-complex}{Lecture 46, Definitions 4.1 and 4.2 and Theorem 4.3.}

\begin{definition}[Finite coefficient categories over a point]
\label{def:point-coefficient-categories}
Let \(D^b_{\mathrm{fin}}(A_{E,n})\) be the full subcategory of
\(D(A_{E,n})\) consisting of bounded complexes with finite cohomology
modules.  Let \(D^b_{\mathrm{fg}}(\mathcal O_E)\) be the full subcategory of
\(D(\mathcal O_E)\) consisting of bounded derived-complete complexes with
finitely generated cohomology modules, and let \(D^b_{\mathrm{fd}}(E)\) be
the bounded derived category of finite-dimensional \(E\)-vector spaces.  Put
\[
 D^b_{\mathrm{fd}}(\Qlb)=
 \Kar\left(\underset{E\in\mathcal E_\ell}{\operatorname{2-colim}}
 D^b_{\mathrm{fd}}(E)\right).
\]
\uses{def:weil-derived-complete-integral,def:qbar-finext-diagram,
  def:filtered-two-colimit,def:qbar-coefficient-category}
\end{definition}

\BlueprintLocalDeclarationRegistry{Definition}{def:point-coefficient-categories}{def:weil-derived-complete-integral,def:qbar-finext-diagram,def:filtered-two-colimit,def:qbar-coefficient-category}{Lecture 36, Definitions 1.1, 1.5, and 1.6.}

\begin{definition}[Ordinary and compactly supported cohomology]
\label{def:cohomology-complexes}
Let \(Y\) be separated and of finite type over \(k\).  For
\(s=\Spec k\), evaluation at its geometric point gives equivalences
\[
\begin{aligned}
 \operatorname{Fib}_s:D^b_c(s,A_{E,n})&\xrightarrow{\sim}
 D^b_{\mathrm{fin}}(A_{E,n}),\\
 \operatorname{Fib}_s:D^b_c(s,\mathcal O_E)&\xrightarrow{\sim}
 D^b_{\mathrm{fg}}(\mathcal O_E),\\
 \operatorname{Fib}_s:D^b_c(s,E)&\xrightarrow{\sim}D^b_{\mathrm{fd}}(E).
\end{aligned}
\]
For \(\Qlb\), the target is \(D^b_{\mathrm{fd}}(\Qlb)\).  For
\(K\in D^b_c(Y,\Lambda)\), define the coefficient complexes
\[
 \RGamma(Y,K)=\operatorname{Fib}_s(Ra_{Y*}K),
 \qquad
 \RGammac(Y,K)=\operatorname{Fib}_s(Ra_{Y!}K),
\]
and
\[
 H^i(Y,K)=\mathrm H^i\bigl(\RGamma(Y,K)\bigr),
 \qquad
 H_c^i(Y,K)=\mathrm H^i\bigl(\RGammac(Y,K)\bigr).
\]
If \(K\) is a Weil complex, each of these four symbols abbreviates the same
construction applied to \(\operatorname{For}(K)\).
Under the finiteness proposition below, these coefficient spaces are
finite-dimensional for field coefficients; their dimensions, traces, and
determinants have their ordinary linear-algebraic meanings.
If \(Y\) is proper, \(Ra_{Y!}\cong Ra_{Y*}\), so
\(H_c^i(Y,K)\cong H^i(Y,K)\).
\uses{def:point-coefficient-categories,def:derived-direct-image,
  def:all-coefficient-compact-support,def:strict-weil-complex,
  prop:constructible-truncations}
\end{definition}

\BlueprintLocalDeclarationRegistry{Definition}{def:cohomology-complexes}{def:point-coefficient-categories,def:derived-direct-image,def:all-coefficient-compact-support,def:strict-weil-complex,prop:constructible-truncations}{Lecture 36, Definitions 2.1 and 2.2 and Theorems 2.3--2.6.}

\begin{definition}[Global geometric Frobenius]
\label{def:global-geometric-frobenius}
For a Weil complex \((K,\theta)\), the geometric Frobenius endomorphism of
ordinary cohomology is induced by
\[
 \RGamma(Y,\operatorname{For}K)\xrightarrow{\sim}
 \RGamma(Y,Lg_Y^*\operatorname{For}K)
 \xrightarrow{\RGamma(\theta_K)}\RGamma(Y,\operatorname{For}K).
\]
The compactly supported operator is defined analogously using \(Ra_{Y!}\).
Both operators, and their actions on cohomology, are denoted \(F\).
\uses{def:weil-frobenius-convention,def:strict-weil-complex,
  def:cohomology-complexes,prop:weil-operations}
\end{definition}

\BlueprintLocalDeclarationRegistry{Definition}{def:global-geometric-frobenius}{def:weil-frobenius-convention,def:strict-weil-complex,def:cohomology-complexes,prop:weil-operations}{Lecture 34, Definition 3.3.}

\begin{definition}[Local Frobenius and local alternating trace]
\label{def:local-frobenius-trace}
Let \(y\in|Y_0|\) be a closed point of degree
\(r=[\kappa(y):\Fq]\).  Choose an embedding
\(\kappa(y)\hookrightarrow k\) and the associated geometric point
\(\bar y\to Y\).  Iterating the Weil isomorphism \(r\) times and using the
resulting return to \(\bar y\) defines the geometric Frobenius operator
\(F_y\) on every stalk
\(\mathcal H^j(\operatorname{For}K)_{\bar y}\); its conjugacy class is
independent of the chosen embedding.

More generally, for \(x\in Y_0(\mathbf F_{q^n})\), the \(n\)-fold iterate
at the associated geometric point defines \(F_x^{(n)}\).  Put
\[
 \Tr(F_x^{(n)}\mid K_{\bar x})=
 \sum_j(-1)^j
 \Tr\bigl(F_x^{(n)}\mid
 \mathcal H^j(\operatorname{For}K)_{\bar x}\bigr).
\]
\uses{def:weil-base-field,def:strict-weil-complex,
  prop:constructible-truncations,def:weil-frobenius-convention}
\end{definition}

\BlueprintLocalDeclarationRegistry{Definition}{def:local-frobenius-trace}{def:weil-base-field,def:strict-weil-complex,prop:constructible-truncations,def:weil-frobenius-convention}{Lecture 34, Definition 3.4; Lecture 46, Definition 2.5.}

\begin{proposition}[Compact-support finiteness and Frobenius]
\label{prop:cohomological-finiteness}
\uses{prop:compact-support-coefficient-passage,prop:weil-operations,
  def:cohomology-complexes,def:global-geometric-frobenius}
Let \(Y_0/\Fq\) be separated and of finite type, let \(Y\) be its geometric
fiber, and let \(K\in D^b_{c,\mathrm W}(Y,E)\), with the same assertion for
\(\Qlb\) after choosing a finite field of definition.  Then
\(H_c^i(Y,K)\), meaning \(H_c^i(Y,\operatorname{For}K)\), is
finite-dimensional, vanishes outside a finite interval,
and carries a geometric Frobenius automorphism.  If \(Y_0\) is proper, the
canonical isomorphism
\[
 H_c^i(Y,K)\xrightarrow{\sim}H^i(Y,K)
\]
is Frobenius-equivariant.
If \(Y\) has dimension \(e\) and \(K=E_Y\) is the constant Weil sheaf, then
\[
 H_c^i(Y,E)=0\qquad\text{for }i<0\text{ or }i>2e.
\]
For proper \(Y\), the same range holds for ordinary cohomology.  In
particular, if \(Y\) is proper of pure dimension \(e\), then
\(H^i(Y,E)=0\) for \(i\notin[0,2e]\).
\end{proposition}

\sketchslot{}


\begin{definition}[Smooth-proper datum and its cohomological family]
\label{def:smooth-proper-cohomological-family}
A \emph{smooth-proper datum} \(\mathscr X\) consists of a smooth, proper,
geometrically connected \(\Fq\)-scheme \(X_0\), an integer \(d\geq0\), an
identification of \(d\) with the pure dimension of \(X_0\), and the
geometric fiber \(X=X_0\times_{\Fq}k\).  For every prime
\(\ell\neq p\), define
\[
 V_{i,\ell}=H^i(X,\Ql),
 \qquad
 F_{i,\ell}=F\vert_{V_{i,\ell}}
 \qquad(0\leq i\leq2d).
\]
Proposition~\ref{prop:cohomological-finiteness} makes \(V_{i,\ell}\) a
finite-dimensional \(\Ql\)-vector space and \(F_{i,\ell}\) an
automorphism, and it gives \(H^i(X,\Ql)=0\) outside the displayed range.
Define
\[
 P_{i,\ell}(T)=
 \det(1-TF_{i,\ell}\mid V_{i,\ell})\in\Ql[T].
\]
It has constant coefficient \(1\), degree
\(\dim_{\Ql}V_{i,\ell}\), and leading coefficient
\((-1)^{\dim V_{i,\ell}}\det(F_{i,\ell})\).
Every displayed symbol is a function of \(\mathscr X\), even when the
argument is suppressed in mathematical notation.
\uses{def:weil-base-field,def:weil-geometric-hypotheses,
  def:cohomology-complexes,def:global-geometric-frobenius,
  prop:cohomological-finiteness}
\end{definition}

\BlueprintLocalDeclarationRegistry{Definition}{def:smooth-proper-cohomological-family}{def:weil-base-field,def:weil-geometric-hypotheses,def:cohomology-complexes,def:global-geometric-frobenius,prop:cohomological-finiteness}{Exact local declaration (Definition~\ref{def:smooth-proper-cohomological-family}).}

\begin{convention}[Fixed smooth-proper datum]
\label{conv:fixed-smooth-proper-datum}
In Sections 5--9 fix one smooth-proper datum \(\mathscr X\).  The symbols
\(X_0,d,X,V_{i,\ell},F_{i,\ell},P_{i,\ell}\) always mean the components or
functions attached to this datum by
Definition~\ref{def:smooth-proper-cohomological-family}.  Every declaration
in these sections has \(\mathscr X\) as an explicit first parameter; the
suppressed argument is only a typesetting convention.
\uses{def:smooth-proper-cohomological-family}
\end{convention}

\BlueprintLocalDeclarationRegistry{Convention}{conv:fixed-smooth-proper-datum}{def:smooth-proper-cohomological-family}{Exact local convention (Convention~\ref{conv:fixed-smooth-proper-datum}).}

\begin{definition}[Ordinary roots and reciprocal roots]
\label{def:ordinary-reciprocal-roots}
Let \(P(T)\) be a polynomial with \(P(0)=1\).  An \emph{ordinary root} is a
scalar \(\beta\) satisfying \(P(\beta)=0\).  It is nonzero.  A
\emph{reciprocal root} is a scalar \(\alpha\) satisfying
\(P(\alpha^{-1})=0\).  Over an algebraic closure one has
\[
 P(T)=\prod_j(1-\alpha_jT),
\]
where the \(\alpha_j\), counted with algebraic multiplicity, are the
reciprocal roots.  For \(P_{i,\ell}\), they are exactly the eigenvalues of
\(F_{i,\ell}\) after scalar extension to \(\Qlb\).
\uses{conv:fixed-smooth-proper-datum}
\end{definition}

\BlueprintLocalDeclarationRegistry{Definition}{def:ordinary-reciprocal-roots}{conv:fixed-smooth-proper-datum}{Lecture 49, Definitions 5.1 and 8.13.}

\section{Point counting and the trace-to-determinant step}

\begin{definition}[Point counts and zeta function]
\label{def:point-count-zeta}
For every finite-type \(\Fq\)-scheme \(Y_0\) and \(n\geq1\), put
\[
 N_n(Y_0)=\#Y_0(\mathbf F_{q^n}).
\]
For a closed point \(x\in|Y_0|\), set
\(\deg(x)=[\kappa(x):\Fq]\).  Define the Hasse--Weil zeta series by
\[
 Z(Y_0,T)=
 \exp\left(\sum_{n\geq1}N_n(Y_0)\frac{T^n}{n}\right)
 \in1+T\mathbf Q[[T]].
\]
The exponential is the formal power-series exponential.
\uses{def:weil-base-field}
\end{definition}

\BlueprintLocalDeclarationRegistry{Definition}{def:point-count-zeta}{def:weil-base-field}{Lecture 49, Definitions 4.1, 4.2, and 4.4.}

\begin{proposition}[Closed-point and integral-series identity]
\label{prop:euler-product}
\uses{def:point-count-zeta}
For every finite-type \(\Fq\)-scheme \(Y_0\), the closed points of degree
dividing \(n\) and the geometric Frobenius orbits give
\[
 \#Y_0(\mathbf F_{q^n})
 =\sum_{r\mid n}r\,c_r,
\]
where \(c_r\) is the number of closed points of degree \(r\).  Consequently
\[
 Z(Y_0,T)
 =\prod_{y\in|Y_0|}(1-T^{\deg(y)})^{-1}
 \in1+T\mathbf Z[[T]].
\]
The product is coefficientwise locally finite: modulo \(T^{m+1}\), only
closed points of degree at most \(m\) occur.
\end{proposition}

\sketchslot{}


\begin{proposition}[Every-iterate Grothendieck--Lefschetz identity]
\label{prop:every-iterate-trace}
\uses{def:point-count-zeta,def:local-frobenius-trace,
  def:global-geometric-frobenius,prop:cohomological-finiteness,
  conv:fixed-smooth-proper-datum}
Let \(Y_0/\Fq\) be separated and of finite type, let \(Y\) be its geometric
fiber, and let \(K\in D^b_{c,\mathrm W}(Y,E)\).  For every \(n\geq1\),
\[
 \sum_{x\in Y_0(\mathbf F_{q^n})}
 \Tr(F_x^{(n)}\mid K_{\bar x})
 =\sum_i(-1)^i
 \Tr\bigl(F^n\mid H_c^i(Y,K)\bigr).
\]
Both sums are finite.  For the constant Weil sheaf \(E_Y\), each local
alternating trace is \(1\), so
\[
 N_n(Y_0)=\sum_i(-1)^i
 \Tr\bigl(F^n\mid H_c^i(Y,E)\bigr).
\]
If \(Y_0\) is proper, compactly supported cohomology may be replaced by
ordinary cohomology.  In particular, for the fixed smooth proper \(X_0\),
\[
 N_n(X_0)=\sum_{i=0}^{2d}(-1)^i
 \Tr(F_{i,\ell}^n\mid V_{i,\ell}).
\]
\end{proposition}

\sketchslot{}


\begin{proposition}[Trace-to-determinant factorization]
\label{prop:trace-determinant}
\uses{prop:every-iterate-trace,def:smooth-proper-cohomological-family,
  def:point-count-zeta}
Let \((V_i,u_i)\) be a finite graded family of finite-dimensional vector
spaces and endomorphisms over a characteristic-zero field.  Put
\[
 P_i(T)=\det(1-Tu_i\mid V_i),
 \qquad
 \Delta(T)=\prod_iP_i(T)^{(-1)^{i+1}}.
\]
Then
\[
 -\log P_i(T)=
 \sum_{n\geq1}\Tr(u_i^n\mid V_i)\frac{T^n}{n},
\]
and, with the indicated sign convention,
\[
 \log\Delta(T)=
 \sum_{n\geq1}\left(\sum_i(-1)^i\Tr(u_i^n\mid V_i)\right)
 \frac{T^n}{n}.
\]
For \(V_i=H_c^i(X,\Ql)\) and \(u_i=F_{i,\ell}\), the every-iterate trace
identity is equivalent, after coefficient extension
\(\mathbf Q[[T]]\to\Ql[[T]]\), to
\[
 Z(X_0,T)=
 \prod_{i=0}^{2d}P_{i,\ell}(T)^{(-1)^{i+1}}.
\]
This identity belongs to \(\Ql(T)\cap\Ql[[T]]\).
\end{proposition}

\sketchslot{}


\section{Weights before arithmetic recovery}

\begin{definition}[Punctual weights]
\label{def:punctual-weights}
Fix an abstract field isomorphism \(\iota:\Qlb\xrightarrow{\sim}\mathbf C\).
Let \(\mathcal F\) be a constructible \(\Qlb\)-Weil sheaf on the geometric
fiber of \(Y_0/\Fq\).  It is \emph{punctually \(\iota\)-pure of weight
\(w\in\mathbf R\)} if, for every closed point \(y\), every eigenvalue
\(\alpha\) of geometric Frobenius on \(\mathcal F_{\bar y}\) satisfies
\[
 |\iota(\alpha)|=q_y^{w/2},
 \qquad q_y=\#\kappa(y).
\]
It is \emph{\(\iota\)-mixed of weights at most \(w\)} if it has a finite
filtration by constructible Weil subsheaves whose graded pieces are
punctually \(\iota\)-pure of respective weights \(w_\nu\) with
\(w_\nu\leq w\).
\uses{def:strict-weil-complex,def:local-frobenius-trace}
\end{definition}

\BlueprintLocalDeclarationRegistry{Definition}{def:punctual-weights}{def:strict-weil-complex,def:local-frobenius-trace}{Lecture 46, Definitions 7.1--7.3; Lecture 47, Definition 2.1.}

\begin{definition}[Weights of a Weil complex]
\label{def:weights-of-complex}
Assume that the Verdier duality functor \(D_Y\) is available on
\(D^b_{ctf}(Y,\Qlb)=D^b_c(Y,\Qlb)\), as in
Proposition~\ref{prop:functor-package}.  For
\(K\in D^b_{c,\mathrm W}(Y,\Qlb)\), write
\(\Wt_{\iota,\leq w}(K)\) if every
\(\mathcal H^j(\operatorname{For}K)\), with its induced Weil structure, is
\(\iota\)-mixed of weights at most \(w+j\).  Write
\(\Wt_{\iota,\geq w}(K)\) if
\(\Wt_{\iota,\leq-w}(D_YK)\) holds.  The complex is
\emph{\(\iota\)-pure of weight \(w\)} if both bounds hold.
\uses{def:punctual-weights,def:strict-weil-complex,
  prop:constructible-truncations,prop:functor-package}
\end{definition}

\BlueprintLocalDeclarationRegistry{Definition}{def:weights-of-complex}{def:punctual-weights,def:strict-weil-complex,prop:constructible-truncations,prop:functor-package}{Lecture 46, Definitions 8.1 and 8.2; Lecture 47, Definition 2.2.}

\begin{definition}[Weil numbers]
\label{def:weil-numbers}
For a number field \(L/\mathbf Q\), define
\[
 \operatorname{Emb}_{\mathbf Q}(L,\mathbf C)=
 \{\tau:L\longrightarrow\mathbf C:
   \tau\text{ is a field homomorphism and }\tau|_{\mathbf Q}
   =\operatorname{id}_{\mathbf Q}\}.
\]
For a nonzero algebraic number \(\alpha\in\overline{\mathbf Q}\), let
\(\mathbf Q(\alpha)\) be the smallest subfield of
\(\overline{\mathbf Q}\) containing \(\mathbf Q\) and \(\alpha\).  The
number \(\alpha\) is a \emph{\(q\)-Weil number of weight \(w\)} if
\[
 \forall\tau\in
 \operatorname{Emb}_{\mathbf Q}(\mathbf Q(\alpha),\mathbf C),
 \qquad |\tau(\alpha)|=q^{w/2}.
\]
It is an \emph{algebraic integer} if there are an integer \(m\geq1\) and
integers \(a_0,\ldots,a_{m-1}\) such that
\[
 \alpha^m+a_{m-1}\alpha^{m-1}+\cdots+a_1\alpha+a_0=0.
\]
\uses{def:weil-base-field}
\end{definition}

\BlueprintLocalDeclarationRegistry{Definition}{def:weil-numbers}{def:weil-base-field}{Lecture 46, Definition 6.1.}

\begin{proposition}[Smooth-proper weight purity]
\label{prop:weight-purity}
\uses{conv:fixed-smooth-proper-datum,def:weights-of-complex,
  def:ordinary-reciprocal-roots,prop:functor-package}
For every abstract field isomorphism
\(\iota:\Qlb\xrightarrow{\sim}\mathbf C\), the following statements hold.
\begin{enumerate}
\item[\textup{(a)}] The constant Weil sheaf \((\Qlb)_X\) is \(\iota\)-pure of weight \(0\).
\item[\textup{(b)}] The Frobenius module \(H_c^i(X,\Qlb)\) is \(\iota\)-mixed of weights
at most \(i\).
\item[\textup{(c)}] There is a Frobenius-equivariant isomorphism
\(H_c^i(X,\Qlb)\xrightarrow{\sim}H^i(X,\Qlb)\).
\item[\textup{(d)}] There is a Frobenius-equivariant perfect duality
\[
 H^i(X,\Qlb)^\vee\xrightarrow{\sim}
 H^{2d-i}(X,\Qlb)(d).
\]
\item[\textup{(e)}] The Frobenius module \(H^i(X,\Qlb)\) is \(\iota\)-pure of weight \(i\).
\item[\textup{(f)}] Every reciprocal root \(\alpha\in\Qlb\) of \(P_{i,\ell}\) satisfies
\[
 |\iota(\alpha)|=q^{i/2}.
\]
\end{enumerate}
\end{proposition}

\sketchslot{}


\begin{definition}[Parity products at a finite realization]
\label{def:realization-parity-products}
For a prime \(\ell\neq p\), define
\[
 N_\ell(T)=\prod_{\substack{0\leq i\leq2d\\ i\text{ odd}}}
 P_{i,\ell}(T),
 \qquad
 D_\ell(T)=\prod_{\substack{0\leq i\leq2d\\ i\text{ even}}}
 P_{i,\ell}(T).
\]
Both products are functions of the fixed smooth-proper datum
\(\mathscr X\).
\uses{conv:fixed-smooth-proper-datum}
\end{definition}

\BlueprintLocalDeclarationRegistry{Definition}{def:realization-parity-products}{conv:fixed-smooth-proper-datum}{Exact local declaration; separation properties are Lecture 49, Lemma 10.1 and Corollary 10.2.}

\begin{proposition}[Weight separation]
\label{prop:weight-separation}
\uses{prop:weight-purity,prop:trace-determinant,
  def:realization-parity-products}
Fix \(\ell\neq p\).  If \(i\neq j\), then \(P_{i,\ell}\) and
\(P_{j,\ell}\) have no common ordinary root in \(\Qlb\) and are coprime
in \(\Qlb[T]\).  The parity products of
Definition~\ref{def:realization-parity-products} are coprime, have constant
coefficient \(1\), and satisfy
\[
 Z(X_0,T)=\frac{N_\ell(T)}{D_\ell(T)}
\]
after coefficient extension to \(\Ql[[T]]\).  More precisely, the ordinary
root multiset of each \(P_{i,\ell}\) is exactly the submultiset of roots of
the appropriate parity product whose reciprocal \(\alpha\) satisfies
\[
 |\iota(\alpha)|=q^{i/2}
\]
for every abstract field isomorphism
\(\iota:\Qlb\xrightarrow{\sim}\mathbf C\).
\end{proposition}

\sketchslot{}


\section{Integral factors and realization independence}

\begin{definition}[Normalized reduced integral fraction]
\label{def:normalized-integral-fraction}
A normalized reduced integral presentation of a series
\(Z(T)\in1+T\mathbf Z[[T]]\cap\mathbf Q(T)\) is a pair
\(A,B\in\mathbf Z[T]\) such that
\[
 A(0)=B(0)=1,
 \qquad
 \gcd(A,B)=1\text{ in }\mathbf Q[T],
 \qquad
 Z(T)=\frac{A(T)}{B(T)}
\]
as a Taylor series at \(T=0\).  Such a pair is unique.
\uses{def:point-count-zeta}
\end{definition}

\BlueprintLocalDeclarationRegistry{Definition}{def:normalized-integral-fraction}{def:point-count-zeta}{Lecture 49, Definitions 7.1 and 8.9 and Lemma 8.11.}

\begin{definition}[Hankel matrix of a power series]
\label{def:hankel-matrix}
For \(C(T)=\sum_{n\geq0}c_nT^n\in K[[T]]\), its Hankel matrix is the
infinite matrix
\[
 \mathcal H(C)=(c_{r+s})_{r,s\geq0}.
\]
Its rank is the supremum of the ranks of its finite square submatrices.  The
series \(C\) is rational over \(K\) if and only if this rank is finite;
equivalently, the sequence \((c_n)\) satisfies a linear recurrence over
\(K\) from some index onward.
\end{definition}

\BlueprintLocalDeclarationRegistry{Definition}{def:hankel-matrix}{none}{Lecture 49, Definition 7.3 and Proposition 7.8.}

\begin{proposition}[Rational descent and normalized integral fraction]
\label{prop:rational-descent}
\uses{prop:euler-product,prop:trace-determinant,
  def:normalized-integral-fraction,def:hankel-matrix}
The zeta function has the following properties:
\begin{enumerate}
\item[\textup{(a)}]
\[
 Z(X_0,T)\in\mathbf Q(T).
\]
\item[\textup{(b)}]
\[
 Z(X_0,T)\in1+T\mathbf Z[[T]]\cap\mathbf Q(T).
\]
\item[\textup{(c)}] There is a unique normalized reduced integral presentation
\[
 Z(X_0,T)=\frac{A(T)}{B(T)},
 \qquad A,B\in\mathbf Z[T],
 \qquad A(0)=B(0)=1,
 \qquad \gcd(A,B)=1\text{ in }\mathbf Q[T].
\]
\end{enumerate}
\end{proposition}

\sketchslot{}


\begin{definition}[Canonical numerator and denominator]
\label{def:canonical-integral-fraction}
Let \((A_{\mathscr X},B_{\mathscr X})\) be the unique normalized reduced
integral presentation of \(Z(X_0,T)\) furnished by
Proposition~\ref{prop:rational-descent}.  In Sections 7--9 write these
polynomials as \(A\) and \(B\).  Thus \(A\) and \(B\) are functions of
\(\mathscr X\), not new existentially chosen variables.
\uses{conv:fixed-smooth-proper-datum,def:normalized-integral-fraction,
  prop:rational-descent}
\end{definition}

\BlueprintLocalDeclarationRegistry{Definition}{def:canonical-integral-fraction}{conv:fixed-smooth-proper-datum,def:normalized-integral-fraction,prop:rational-descent}{Exact local canonical choice; existence and uniqueness are Lecture 49, Theorem 8.12.}

\begin{definition}[Integral realization of a Frobenius factor]
\label{def:integral-realization-factor}
For \(P(T)=\sum_ra_rT^r\in\mathbf Z[T]\), write
\[
 P^{(\ell)}(T)=\sum_r(a_r\cdot1_{\Ql})T^r\in\Ql[T].
\]
An integral realization-independent factor in degree \(i\) is a polynomial
\(P_i\in\mathbf Z[T]\) such that
\[
 P_i^{(\ell)}=P_{i,\ell}
\]
for every prime \(\ell\neq p\).
\uses{conv:fixed-smooth-proper-datum}
\end{definition}

\BlueprintLocalDeclarationRegistry{Definition}{def:integral-realization-factor}{conv:fixed-smooth-proper-datum}{Exact local declaration; realization-family input is Lecture 49, Definition 12.1.}

\begin{definition}[The algebraic closure of \(\mathbf Q\) inside
\(\Qlb\)]
\label{def:algebraic-numbers-in-qbarl}
Set
\[
 \overline{\mathbf Q}^{(\ell)}=
 \{a\in\Qlb:a\text{ is algebraic over }\mathbf Q\}.
\]
This is an algebraically closed subfield of \(\Qlb\).  After choosing an
algebraic closure \(\overline{\mathbf Q}\) of \(\mathbf Q\), a
\(\mathbf Q\)-embedding
\(\eta:\overline{\mathbf Q}\xrightarrow{\sim}
\overline{\mathbf Q}^{(\ell)}\) is called an ell-adic realization of
algebraic numbers.
\uses{def:weil-base-field}
\end{definition}

\BlueprintLocalDeclarationRegistry{Definition}{def:algebraic-numbers-in-qbarl}{def:weil-base-field}{Lecture 49, Definition 9.1 and Proposition 9.2.}

\begin{lemma}[Extension of embeddings and intrinsic weights]
\label{lem:embedding-extension}
\uses{def:algebraic-numbers-in-qbarl,def:weil-numbers}
There exists a \(\mathbf Q\)-isomorphism
\[
 \eta:\overline{\mathbf Q}\xrightarrow{\sim}
 \overline{\mathbf Q}^{(\ell)}.
\]
Let \(L\subset\Qlb\) be a number field and let
\(\sigma:L\hookrightarrow\mathbf C\) be a field embedding.  There is an
abstract field isomorphism
\[
 \iota:\Qlb\xrightarrow{\sim}\mathbf C
\]
whose restriction to \(L\) is \(\sigma\).  If
\(\alpha\in\overline{\mathbf Q}\) satisfies
\[
 |\iota(\eta(\alpha))|=q^{i/2}
\]
for every ell-adic realization
\(\eta:\overline{\mathbf Q}\xrightarrow{\sim}
\overline{\mathbf Q}^{(\ell)}\) and every abstract field isomorphism
\(\iota:\Qlb\xrightarrow{\sim}\mathbf C\), then
\[
 |\tau(\alpha)|=q^{i/2}
\]
for every \(\tau:\mathbf Q(\alpha)\hookrightarrow\mathbf C\).
The set of \(q\)-Weil numbers of weight \(i\) is stable under
\(\operatorname{Gal}(\overline{\mathbf Q}/\mathbf Q)\).
\end{lemma}

\sketchslot{}


\begin{definition}[Weight-stratum multisets and polynomials]
\label{def:weight-stratum-factors}
Let \(A,B\) be the canonical numerator and denominator of
Definition~\ref{def:canonical-integral-fraction}.  For a polynomial
\(C\in\{A,B\}\) of degree \(r\), put
\[
 C^\#(U)=U^rC(U^{-1}).
\]
For odd \(i\), let \(R_i\) be the multiset of roots of \(A^\#\) that are
\(q\)-Weil numbers of weight \(i\); for even \(i\), use the roots of
\(B^\#\).  Multiplicities are inherited from the corresponding reciprocal
polynomial.  Define
\[
 Q_i(T)=\prod_{\alpha\in R_i}(1-\alpha T).
\]
The symbols \(R_i\) and \(Q_i\) are functions of \(\mathscr X\).
\uses{def:canonical-integral-fraction,def:weil-numbers}
\end{definition}

\BlueprintLocalDeclarationRegistry{Definition}{def:weight-stratum-factors}{def:canonical-integral-fraction,def:weil-numbers}{Exact local declaration; Lecture 49, Definition 9.8 and Proposition 9.9.}

\begin{proposition}[Galois-stable weight-stratum factors]
\label{prop:galois-weight-strata}
\uses{def:weight-stratum-factors,prop:weight-separation,
  lem:embedding-extension}
The reciprocal polynomials \(A^\#\) and \(B^\#\) are monic, and their
roots, with multiplicity, are algebraic integers and are the reciprocal
roots of \(A\) and \(B\).  The multisets and polynomials of
Definition~\ref{def:weight-stratum-factors} satisfy:
\begin{enumerate}
\item[\textup{(a)}] each \(R_i\) is stable, with multiplicities, under
\(\operatorname{Gal}(\overline{\mathbf Q}/\mathbf Q)\);
\item[\textup{(b)}] the normalized stratum polynomial
\[
 Q_i(T)=\prod_{\alpha\in R_i}(1-\alpha T)
\]
has coefficients in \(\mathbf Q\), has algebraic-integer coefficients, and
belongs to \(\mathbf Z[T]\);
\item[\textup{(c)}] for every \(\ell\neq p\) and every ell-adic realization \(\eta\),
scalar extension sends the multiset \(R_i\) to the reciprocal-root multiset
of \(P_{i,\ell}\), and
\[
 Q_i^{(\ell)}(T)=P_{i,\ell}(T).
\]
\end{enumerate}
\end{proposition}

\sketchslot{}


\begin{definition}[Canonical integral cohomological factors]
\label{def:canonical-cohomological-factors}
For \(0\leq i\leq2d\), define
\[
 P_i=P_i(\mathscr X):=Q_i,
\]
where \(Q_i\) is the weight-stratum polynomial of
Definition~\ref{def:weight-stratum-factors}.  Proposition
~\ref{prop:galois-weight-strata} supplies the theorem that
\(P_i\in\mathbf Z[T]\), that \(P_i(0)=1\), and that every finite
realization is \(P_{i,\ell}\).
\uses{def:weight-stratum-factors,prop:galois-weight-strata}
\end{definition}

\BlueprintLocalDeclarationRegistry{Definition}{def:canonical-cohomological-factors}{def:weight-stratum-factors,prop:galois-weight-strata}{Exact local canonical choice; Lecture 49, Theorem 12.3 and Corollary 12.5.}

\begin{proposition}[Integral factor recovery]
\label{prop:integral-recovery}
\uses{def:canonical-integral-fraction,def:realization-parity-products,
  prop:weight-separation,prop:galois-weight-strata,
  def:integral-realization-factor}
There are unique polynomials
\[
 P_i(T)\in\mathbf Z[T]
 \qquad(0\leq i\leq2d)
\]
with \(P_i(0)=1\) and \(P_i^{(\ell)}=P_{i,\ell}\) for every
\(\ell\neq p\).  They equal the polynomials \(Q_i\) of
Proposition~\ref{prop:galois-weight-strata}.  For every \(\ell\neq p\),
\[
 A^{(\ell)}=N_\ell,
 \qquad
 B^{(\ell)}=D_\ell.
\]
The resulting factors satisfy
\[
 Z(X_0,T)=
 \prod_{i=0}^{2d}P_i(T)^{(-1)^{i+1}}
 \quad\text{in }\mathbf Q(T),
\]
and
\[
 \deg P_i=\dim_{\Ql}H^i(X,\Ql)
\]
for every \(\ell\neq p\).
\end{proposition}

\sketchslot{}


\begin{proposition}[Arithmetic roots and realization-independent degrees]
\label{prop:arithmetic-roots}
\uses{def:canonical-cohomological-factors,prop:weight-purity,
  prop:galois-weight-strata,lem:embedding-extension}
Let \(r_i=\deg P_i\).  The reciprocal polynomial
\[
 P_i^\#(U)=U^{r_i}P_i(U^{-1})\in\mathbf Z[U]
\]
is monic, and every reciprocal root \(\alpha\) of \(P_i\) is an algebraic
integer.  For every embedding
\(\tau:\mathbf Q(\alpha)\hookrightarrow\mathbf C\),
\[
 |\tau(\alpha)|=q^{i/2};
\]
and \(\alpha\) is a \(q\)-Weil number of weight \(i\).  For every field
embedding \(\eta:\overline{\mathbf Q}\hookrightarrow\Qlb\), the scalar
\(\eta(\alpha)\) occurs as an eigenvalue of geometric Frobenius on
\(H^i(X,\Qlb)\).  Every ordinary root \(\beta=\alpha^{-1}\) satisfies,
for every \(\tau:\mathbf Q(\beta)\hookrightarrow\mathbf C\),
\[
 |\tau(\beta)|=q^{-i/2}.
\]
Finally,
\[
 r_i=\dim_{\Ql}H^i(X,\Ql)
\]
for every \(\ell\neq p\); both the factor and its degree are independent of
the chosen ell-adic realization.
\end{proposition}

\sketchslot{}


\begin{definition}[Betti numbers]
\label{def:betti-numbers}
For \(0\leq i\leq2d\), define
\[
 b_i=\deg P_i.
\]
For \(i\notin[0,2d]\), put \(b_i=0\).  Thus \(b_i\) is defined solely from
the canonical integral factor and is independent of a realization by
definition; its cohomological interpretation is
Proposition~\ref{prop:arithmetic-roots}.
\uses{def:canonical-cohomological-factors}
\end{definition}

\BlueprintLocalDeclarationRegistry{Definition}{def:betti-numbers}{def:canonical-cohomological-factors}{Exact local declaration; independence is Lecture 49, Corollary 12.4.}

\section{Duality, endpoints, reciprocity, and the functional equation}

\begin{theorem}[Exceptional admissibility of a smooth-proper structure map]
\label{thm:smooth-proper-structure-admissible}
\uses{def:exceptional-adjoint-admissible,conv:fixed-smooth-proper-datum}
For the fixed smooth-proper datum, the structure morphism
\(a_X:X\to\Spec k\) is exceptional-adjoint admissible at every finite
coefficient stage, at \(\mathcal O_E\), at \(E\), and at \(\Qlb\), with the
uniform amplitude and coefficient-comparison data required by
Definition~\ref{def:exceptional-adjoint-admissible}.
\sorryleaf{Lecture 44, Theorem 6.20, together with the smooth-proper
amplitude and finite-Tor admissibility input made explicit there.}
\end{theorem}

\sketchslot{}

\begin{convention}[Exceptional admissibility for the structure morphism]
\label{conv:structure-map-admissibility}
Throughout this section, the smooth proper structure morphism
\(a_X:X\to\Spec k\) is equipped with the admissibility datum of
Theorem~\ref{thm:smooth-proper-structure-admissible}.
\uses{thm:smooth-proper-structure-admissible}
\end{convention}

\BlueprintLocalDeclarationRegistry{Convention}{conv:structure-map-admissibility}{thm:smooth-proper-structure-admissible}{Exact local convention, discharged by Theorem~\ref{thm:smooth-proper-structure-admissible}.}

\begin{definition}[Poincare trace and cup-product pairing]
\label{def:poincare-trace-pairing}
For the smooth structure morphism \(a_X:X\to\Spec k\), smooth purity gives
\[
 Ra_X^!\Ql\cong\Ql(d)[2d].
\]
Through this isomorphism, the counit
\(Ra_{X!}Ra_X^!\Ql\to\Ql\) defines the trace map
\[
 \operatorname{tr}_X:H_c^{2d}(X,\Ql(d))\longrightarrow\Ql.
\]
Properness identifies compactly supported and ordinary cohomology.  After
untwisting, write
\[
 \operatorname{tr}^{(-d)}_X:H^{2d}(X,\Ql)\longrightarrow\Ql(-d).
\]
The derived tensor product and multiplication
\((\Ql)_X\otimes^L(\Ql)_X\to(\Ql)_X\) define cup products
\[
 \smile:H^i(X,\Ql)\otimes_{\Ql}H^j(X,\Ql)
 \longrightarrow H^{i+j}(X,\Ql).
\]
For every \(i\), define
\[
 \langle x,y\rangle_i=
 \operatorname{tr}^{(-d)}_X(x\smile y):
 H^i(X,\Ql)\times H^{2d-i}(X,\Ql)
 \longrightarrow\Ql(-d).
\]
\uses{conv:fixed-smooth-proper-datum,conv:structure-map-admissibility,
  def:exceptional-inverse-image,def:coefficient-twists,
  def:derived-pullback-tensor-hom,def:cohomology-complexes}
\end{definition}

\BlueprintLocalDeclarationRegistry{Definition}{def:poincare-trace-pairing}{conv:fixed-smooth-proper-datum,conv:structure-map-admissibility,def:exceptional-inverse-image,def:coefficient-twists,def:derived-pullback-tensor-hom,def:cohomology-complexes}{Lecture 36, Definitions 4.3 and 4.5.}

\begin{proposition}[Poincare duality and Frobenius similitude]
\label{prop:poincare-pairing}
\uses{def:poincare-trace-pairing,def:frobenius-on-twists,
  prop:weil-operations}
The trace map is an isomorphism, and for every \(i\) the cup-product pairing
\(\langle-,-\rangle_i\) is perfect.  Equivalently, it induces a
Frobenius-equivariant isomorphism
\[
 H^{2d-i}(X,\Ql)
 \xrightarrow{\sim}H^i(X,\Ql)^\vee(-d).
\]
Since geometric Frobenius acts on \(\Ql(-d)\) by \(q^d\),
\[
 \langle Fx,Fy\rangle_i=q^d\langle x,y\rangle_i.
\]
Graded commutativity gives
\[
 \langle y,x\rangle_{2d-i}=(-1)^i\langle x,y\rangle_i;
\]
in middle degree the pairing is \((-1)^d\)-symmetric.
\end{proposition}

\sketchslot{}


\begin{proposition}[Endpoint cohomology and endpoint factors]
\label{prop:endpoints}
\uses{prop:poincare-pairing,def:canonical-cohomological-factors,
  def:betti-numbers,prop:cohomological-finiteness}
There are Frobenius-equivariant isomorphisms
\[
 H^0(X,\Ql)\cong\Ql,
 \qquad
 H^{2d}(X,\Ql)\cong\Ql(-d).
\]
Geometric Frobenius acts by \(1\) on the first space and by \(q^d\) on the
second.  Consequently
\[
 P_0(T)=1-T,
 \qquad
 P_{2d}(T)=1-q^dT,
 \qquad
 b_0=b_{2d}=1.
\]
\end{proposition}

\sketchslot{}


\begin{definition}[Determinant constants and Euler characteristic]
\label{def:determinant-constants}
For \(0\leq i\leq2d\), set
\[
 D_i=(-1)^{b_i}[T^{b_i}]P_i(T)\in\mathbf Z\setminus\{0\}.
\]
Its image in \(\Ql\) is
\(\det(F_{i,\ell}\mid H^i(X,\Ql))\).  Define
\[
 \chi=\sum_{i=0}^{2d}(-1)^ib_i.
\]
Define also the alternating determinant scalar
\[
 \Delta=\prod_{i=0}^{2d}D_i^{(-1)^i}\in\mathbf Q^\times.
\]
\uses{def:canonical-cohomological-factors,def:betti-numbers,
  conv:fixed-smooth-proper-datum}
\end{definition}

\BlueprintLocalDeclarationRegistry{Definition}{def:determinant-constants}{def:canonical-cohomological-factors,def:betti-numbers,conv:fixed-smooth-proper-datum}{Exact local packaging; common determinants are Lecture 49, Definition 17.1.}

\begin{proposition}[Duality reciprocity]
\label{prop:duality-reciprocity}
\uses{prop:poincare-pairing,def:canonical-cohomological-factors,
  def:betti-numbers,def:determinant-constants}
For every \(0\leq i\leq2d\), Poincare duality identifies the multiset of
Frobenius eigenvalues in degree \(2d-i\) with
\[
 \left\{\frac{q^d}{\alpha}:\alpha
 \text{ is an eigenvalue in degree }i\right\},
\]
including algebraic multiplicity.  Moreover,
\[
 b_i=b_{2d-i},
 \qquad
 D_iD_{2d-i}=q^{db_i},
\]
and the integral factors satisfy the Laurent-polynomial identity
\[
 P_{2d-i}(q^{-d}T^{-1})
 =(-1)^{b_i}T^{-b_i}D_i^{-1}P_i(T)
 \quad\text{in }\mathbf Q[T,T^{-1}].
\]
For \(i=d\), the middle pairing is \((-1)^d\)-symmetric.  In particular,
if \(d\) is odd, then \(b_d\) is even.
\end{proposition}

\sketchslot{}


\begin{definition}[Normalized functional-equation sign]
\label{def:functional-equation-sign}
Duality reciprocity implies that \(d\chi\) is even.  Define
\[
 h=\frac{d\chi}{2}\in\mathbf Z
\]
and
\[
 \varepsilon=
 (-1)^\chi q^{-h}\Delta
 \in\mathbf Q^\times.
\]
Negative exponents in \(\Delta\) denote multiplicative inverses in
\(\mathbf Q^\times\).
\uses{def:determinant-constants,prop:duality-reciprocity}
\end{definition}

\BlueprintLocalDeclarationRegistry{Definition}{def:functional-equation-sign}{def:determinant-constants,prop:duality-reciprocity}{Lecture 49, Definition 17.4 and Theorem 17.6.}

\begin{proposition}[Functional equation]
\label{prop:functional-equation}
\uses{prop:trace-determinant,prop:integral-recovery,
  prop:duality-reciprocity,def:functional-equation-sign}
The alternating determinant scalar satisfies
\[
 \Delta^2=q^{d\chi}.
\]
The normalized sign satisfies
\[
 \varepsilon^2=1,\qquad \varepsilon\in\{1,-1\},
\]
and the zeta function satisfies
\[
 Z(X_0,q^{-d}T^{-1})
 =\varepsilon q^{d\chi/2}T^\chi Z(X_0,T)
 \quad\text{in }\mathbf Q(T).
\]
The same identity in terms of the determinant constants is
\[
 Z(X_0,q^{-d}T^{-1})
 =(-1)^\chi T^\chi
 \left(\prod_{i=0}^{2d}D_i^{(-1)^i}\right)Z(X_0,T).
\]
\end{proposition}

\sketchslot{}


\section{Optional characteristic-zero comparison}

\begin{definition}[Smooth proper lift]
\label{def:smooth-proper-lift}
A smooth proper lift of \(X_0\) consists of a mixed-characteristic discrete
valuation ring \(R\) with residue field \(\Fq\), fraction field \(K\), and a
smooth proper \(R\)-scheme \(\mathcal X\) together with an identification
\[
 \mathcal X\times_R\Fq\cong X_0.
\]
A lift datum also includes an embedding \(K\hookrightarrow\mathbf C\), an
algebraic closure \(\overline K\), an extension
\(\overline K\hookrightarrow\mathbf C\), and compatible geometric generic
and special points of \(\Spec R\).  For such a datum \(\mathscr L\), put
\[
 X_{\mathbf C}(\mathscr L)=\mathcal X_K\times_K\mathbf C.
\]
\uses{conv:fixed-smooth-proper-datum}
\end{definition}

\BlueprintLocalDeclarationRegistry{Definition}{def:smooth-proper-lift}{conv:fixed-smooth-proper-datum}{Exact local declaration (Definition~\ref{def:smooth-proper-lift}).}

\begin{proposition}[Betti comparison for a supplied lift]
\label{prop:betti-comparison}
\uses{def:smooth-proper-lift,def:betti-numbers,
  prop:coefficient-passage}
For every smooth proper lift datum \(\mathscr L\) as above and every \(i\),
smooth proper base
change and comparison with complex topology give isomorphisms
\[
 H^i(X,\Ql)
 \cong H^i_{\mathrm{et}}(\mathcal X_{\overline K},\Ql)
 \cong H^i_{\mathrm{sing}}
 \bigl(X_{\mathbf C}(\mathscr L)(\mathbf C),\mathbf Q\bigr)
 \otimes_{\mathbf Q}\Ql.
\]
Consequently
\[
 b_i=\dim_{\mathbf Q}H^i_{\mathrm{sing}}
 \bigl(X_{\mathbf C}(\mathscr L)(\mathbf C),\mathbf Q\bigr).
\]
This statement is conditional on the existence of the displayed lift; no
lifting hypothesis is needed for the other propositions.
\end{proposition}

\sketchslot{}


\section{Statement-only leaves used by the terminal assembly}

The terminal proof is the only proof expanded in this blueprint.  Each row
below is a statement boundary whose proof is not repeated here.  Its exact
mathematical source and deferred status are recorded explicitly.  The table
includes the principal infrastructure leaves as well as the arithmetic and
duality leaves called by the final proof.

{%
\small
\setlength{\tabcolsep}{3pt}
\setlength\LTleft{\fill}
\setlength\LTright{\fill}
\newcommand{\BlueprintLeafRow}[3]{%
  \hyperref[#1]{\nolinkurl{#1}~(\ref*{#1})} & #2 & \textit{#3}\\}
\begin{longtable}{@{}
  p{0.27\textwidth}
  p{0.58\textwidth}
  p{0.09\textwidth}@{}}
\caption{Statement-only leaves for the terminal assembly.}
\label{tab:terminal-proof-leaves}\\
\toprule
Local label & Exact source & Proof status\\
\midrule
\endfirsthead
\toprule
Local label & Exact source & Proof status\\
\midrule
\endhead
\midrule
\multicolumn{3}{r}{\emph{Continued on next page}}\\
\endfoot
\bottomrule
\endlastfoot
\BlueprintLeafRow{thm:qbar-two-colimit-model}{Lecture 42, Theorems 3.12, 4.5, 5.4, 6.4, and 6.5.}{Deferred}
\BlueprintLeafRow{thm:qbar-functor-descent}{Lecture 42, Definition 8.1 and Theorem 8.2.}{Deferred}
\BlueprintLeafRow{prop:functor-package}{Lecture 40, Theorem 11.1; Lecture 30, Theorem 9.3; Lecture 31, Theorems 9.6 and 12.25; Lecture 33, Theorem 11.24; Lecture 43, Theorem 9.17; Lecture 41, Corollary 7.13; Lecture 42, Theorems 5.4, 6.5, and 10.29; Lecture 44, Theorem 6.20; Lecture 45, Theorem 7.5.}{Deferred}
\BlueprintLeafRow{prop:cohomological-finiteness}{Lecture 30, Theorem 9.3; Lecture 33, Theorem 11.24; Lecture 43, Theorem 9.17.}{Deferred}
\BlueprintLeafRow{prop:euler-product}{Lecture 49, Corollary 4.5.}{Deferred}
\BlueprintLeafRow{prop:every-iterate-trace}{Lecture 34, Theorem 11.5 and Corollary 11.6.}{Deferred}
\BlueprintLeafRow{prop:trace-determinant}{Lecture 49, Theorem 5.6, Proposition 6.3, and Corollary 6.4.}{Deferred}
\BlueprintLeafRow{prop:weight-purity}{Lecture 47, Corollary 10.6; Lecture 35, Theorems 7.3 and 7.8; Lecture 36, Theorem 8.9 and Corollary 8.10.}{Deferred}
\BlueprintLeafRow{prop:weight-separation}{Lecture 49, Lemma 10.1, Corollary 10.2, and Proposition 10.4.}{Deferred}
\BlueprintLeafRow{prop:rational-descent}{Lecture 49, Theorem 7.11, Corollary 7.12, and Theorem 8.12.}{Deferred}
\BlueprintLeafRow{lem:embedding-extension}{Lecture 49, Proposition 9.2 and Lemma 9.7.}{Deferred}
\BlueprintLeafRow{prop:galois-weight-strata}{Lecture 49, Proposition 9.9, Proposition 10.4, and Propositions 11.3--11.4.}{Deferred}
\BlueprintLeafRow{thm:smooth-proper-structure-admissible}{Lecture 44, Theorem 6.20, together with the explicit smooth-proper amplitude and finite-Tor admissibility input.}{Deferred}
\BlueprintLeafRow{prop:poincare-pairing}{Lecture 35, Theorems 7.3 and 7.8; Lecture 36, Theorem 8.9 and Corollary 8.10.}{Deferred}
\BlueprintLeafRow{prop:endpoints}{Exact local leaf from geometric connectedness, cohomological finiteness, and Proposition~\ref{prop:poincare-pairing}.}{Deferred}
\BlueprintLeafRow{prop:duality-reciprocity}{Lecture 49, Theorem 13.5.}{Deferred}
\BlueprintLeafRow{prop:betti-comparison}{Exact local conditional leaf invoking proper base change, smooth-proper specialization, and Artin comparison for the supplied lift.}{Deferred}
\end{longtable}
}%

\section{Terminal theorem}

\begin{theorem}[Weil conjectures for a smooth proper variety]
\label{thm:weil-final}
\uses{def:smooth-proper-cohomological-family,def:point-count-zeta,
  prop:cohomological-finiteness,prop:euler-product,
  prop:every-iterate-trace,prop:trace-determinant,prop:weight-purity,
  prop:weight-separation,prop:rational-descent,
  prop:galois-weight-strata,lem:embedding-extension,
  def:ordinary-reciprocal-roots,def:weil-numbers,
  def:integral-realization-factor,def:realization-parity-products,
  def:canonical-integral-fraction,def:determinant-constants,
  def:canonical-cohomological-factors,def:betti-numbers,
  prop:poincare-pairing,prop:endpoints,prop:duality-reciprocity,
  def:functional-equation-sign,def:smooth-proper-lift,
  prop:betti-comparison}
Let \(X_0\) be a smooth, proper, geometrically connected scheme of pure
dimension \(d\) over \(\Fq\).  Put \(X=X_0\times_{\Fq}k\), let
\(\mathscr X\) be the resulting smooth-proper datum, and let
\(P_i=P_i(\mathscr X)\) be the canonical factors of
Definition~\ref{def:canonical-cohomological-factors}.  Then the following
statements hold.
\begin{enumerate}
\item[\textup{(i)}] The zeta series has both point-counting descriptions
\[
 Z(X_0,T)=
 \exp\left(\sum_{n\geq1}
 \#X_0(\mathbf F_{q^n})\frac{T^n}{n}\right)
 =\prod_{x\in|X_0|}(1-T^{\deg(x)})^{-1},
\]
and belongs to \(1+T\mathbf Z[[T]]\cap\mathbf Q(T)\).

\item[\textup{(ii)}] The polynomials \(P_i\) are the unique polynomials
\[
 P_i(T)\in\mathbf Z[T]
 \qquad(0\leq i\leq2d)
\]
with \(P_i(0)=1\) such that, for every \(\ell\neq p\),
\[
 P_i^{(\ell)}(T)=
 \det(1-TF\mid H^i(X,\Ql)).
\]
They give the rational factorization
\[
 Z(X_0,T)=
 \prod_{i=0}^{2d}P_i(T)^{(-1)^{i+1}}.
\]

\item[\textup{(iii)}] For every \(0\leq i\leq2d\), every reciprocal root \(\alpha\) of
\(P_i\) is an algebraic integer
and a \(q\)-Weil number of weight \(i\):
\[
 |\tau(\alpha)|=q^{i/2}
\]
for every embedding
\(\tau:\mathbf Q(\alpha)\hookrightarrow\mathbf C\).  For every ordinary
root \(\beta\) of \(P_i\) and every embedding
\(\tau:\mathbf Q(\beta)\hookrightarrow\mathbf C\), one has
\(|\tau(\beta)|=q^{-i/2}\).

\item[\textup{(iv)}] For every \(0\leq i\leq2d\), Poincare duality supplies a
Frobenius-equivariant isomorphism
\[
 H^{2d-i}(X,\Ql)\xrightarrow{\sim}H^i(X,\Ql)^\vee(-d).
\]
Consequently, with \(b_i=\deg P_i\),
\[
 b_i=\dim_{\Ql}H^i(X,\Ql),
 \qquad b_i=b_{2d-i},
 \qquad b_0=b_{2d}=1,
\]
independently of \(\ell\).  Moreover
\[
 P_0(T)=1-T,
 \qquad
 P_{2d}(T)=1-q^dT.
\]

\item[\textup{(v)}] For every \(0\leq i\leq2d\), under the Poincare-duality
isomorphism, an eigenvalue \(\alpha\) of geometric Frobenius in degree
\(i\) corresponds, with algebraic multiplicity, to the eigenvalue
\(q^d/\alpha\) in degree \(2d-i\).  If
\[
 D_i=(-1)^{b_i}[T^{b_i}]P_i(T),
 \qquad
 \chi=\sum_{i=0}^{2d}(-1)^ib_i,
\]
then \(d\chi\) is even and, for every \(0\leq i\leq2d\),
\[
 P_{2d-i}(q^{-d}T^{-1})
 =(-1)^{b_i}T^{-b_i}D_i^{-1}P_i(T),
\]
and
\[
 Z(X_0,q^{-d}T^{-1})
 =\varepsilon q^{d\chi/2}T^\chi Z(X_0,T)
\]
for the sign
\[
 \varepsilon=
 (-1)^\chi q^{-d\chi/2}
 \prod_{i=0}^{2d}D_i^{(-1)^i}
 \in\{1,-1\}.
\]

\item[\textup{(vi)}] For every smooth proper lift \(\mathscr L\) of \(X_0\) to
characteristic zero as in Definition~\ref{def:smooth-proper-lift}, with
associated complex fiber \(X_{\mathbf C}(\mathscr L)\), and for every
\(0\leq i\leq2d\),
\[
 b_i=\dim_{\mathbf Q}H^i_{\mathrm{sing}}
 \bigl(X_{\mathbf C}(\mathscr L)(\mathbf C),\mathbf Q\bigr).
\]
\end{enumerate}
\end{theorem}

\sketchslot{Expanded in the terminal proof immediately below.}

\begin{proof}
\uses{def:smooth-proper-cohomological-family,def:point-count-zeta,
  prop:cohomological-finiteness,prop:euler-product,
  prop:every-iterate-trace,prop:trace-determinant,
  prop:weight-purity,prop:weight-separation,prop:rational-descent,
  def:canonical-integral-fraction,prop:galois-weight-strata,
  lem:embedding-extension,def:ordinary-reciprocal-roots,
  def:weil-numbers,def:integral-realization-factor,
  def:realization-parity-products,def:canonical-cohomological-factors,
  def:betti-numbers,prop:poincare-pairing,prop:endpoints,
  def:determinant-constants,
  prop:duality-reciprocity,def:functional-equation-sign,
  def:smooth-proper-lift,prop:betti-comparison}

Fix a prime \(\ell\neq p\).  Proposition
~\ref{prop:cohomological-finiteness} makes
\(V_{i,\ell}=H^i(X,\Ql)\) finite-dimensional, zero for
\(i\notin[0,2d]\), and equipped with the geometric-Frobenius automorphism
\(F_{i,\ell}\).  Definition
~\ref{def:smooth-proper-cohomological-family} therefore gives
\[
 P_{i,\ell}(T)=
 \det(1-TF_{i,\ell}\mid V_{i,\ell})\in\Ql[T].
\]
For every \(n\geq1\), Proposition~\ref{prop:every-iterate-trace} gives
\[
 \#X_0(\mathbf F_{q^n})
 =\sum_{i=0}^{2d}(-1)^i
   \Tr(F_{i,\ell}^n\mid V_{i,\ell}).
\]
Using this equality coefficient by coefficient in \(\Ql[[T]]\) and then
using Proposition~\ref{prop:trace-determinant}, we obtain
\begin{equation}
\begin{aligned}
 \log Z(X_0,T)
 &=\sum_{n\geq1}\#X_0(\mathbf F_{q^n})\frac{T^n}{n}\\
 &=\sum_{i=0}^{2d}(-1)^i
   \sum_{n\geq1}\Tr(F_{i,\ell}^n\mid V_{i,\ell})\frac{T^n}{n}\\
 &=\sum_{i=0}^{2d}(-1)^{i+1}\log P_{i,\ell}(T).
\end{aligned}
\label{eq:weil-final-log-determinant}
\end{equation}
The formal logarithm is injective on \(1+T\Ql[[T]]\), so
\begin{equation}
 Z(X_0,T)=
 \prod_{i=0}^{2d}P_{i,\ell}(T)^{(-1)^{i+1}}
 \quad\text{in }\Ql[[T]].
 \label{eq:weil-final-fixed-ell-factorization}
\end{equation}

\smallskip
\noindent\emph{Clause \textup{(i)}.}
The exponential expression is Definition~\ref{def:point-count-zeta}.
Proposition~\ref{prop:euler-product}, applied to \(X_0\), gives
\[
 Z(X_0,T)=\prod_{x\in|X_0|}(1-T^{\deg(x)})^{-1}
 \quad\text{and}\quad
 Z(X_0,T)\in1+T\mathbf Z[[T]].
\]
Proposition~\ref{prop:rational-descent}, using
Equation~\ref{eq:weil-final-fixed-ell-factorization}, gives
\(Z(X_0,T)\in\mathbf Q(T)\).  This proves \textup{(i)}.

\smallskip
\noindent\emph{Clause \textup{(ii)}.}
Proposition~\ref{prop:weight-purity} says that every reciprocal root of
\(P_{i,\ell}\) has absolute value \(q^{i/2}\) under every abstract
isomorphism \(\Qlb\simeq\mathbf C\).  Proposition
~\ref{prop:weight-separation} therefore gives the coprime normalized
products
\[
 N_\ell=\prod_{\substack{0\leq i\leq2d\\i\text{ odd}}}P_{i,\ell},
 \qquad
 D_\ell=\prod_{\substack{0\leq i\leq2d\\i\text{ even}}}P_{i,\ell},
 \qquad
 Z=\frac{N_\ell}{D_\ell}.
\]
Let \((A,B)\) be the canonical normalized reduced integral fraction of
Definition~\ref{def:canonical-integral-fraction}.  Scalar extension to
\(\Ql[T]\) is again normalized and reduced.  Since
\(\Ql(T)\hookrightarrow\Ql((T))\) is injective, uniqueness of normalized
reduced fractions gives
\begin{equation}
 A^{(\ell)}=N_\ell,
 \qquad
 B^{(\ell)}=D_\ell.
 \label{eq:weil-final-parity-comparison}
\end{equation}
Proposition~\ref{prop:galois-weight-strata} and Definition
~\ref{def:canonical-cohomological-factors} give
\[
 P_i(T)\in\mathbf Z[T],
 \qquad P_i(0)=1,
 \qquad P_i^{(\ell)}(T)=P_{i,\ell}(T)
\]
for every \(\ell\neq p\).  Equation
Equation~\ref{eq:weil-final-parity-comparison}, followed by injectivity of
\(\mathbf Z[T]\to\Ql[T]\), identifies \(A\) with the product of the odd
\(P_i\)'s and \(B\) with the product of the even \(P_i\)'s.  Hence
\[
 Z(X_0,T)=\frac{A(T)}{B(T)}
 =\prod_{i=0}^{2d}P_i(T)^{(-1)^{i+1}}
 \quad\text{in }\mathbf Q(T).
\]
If \(\widetilde P_i\in\mathbf Z[T]\) is another normalized family with
the same realization property, then for one fixed \(\ell\neq p\) its
image equals \(P_{i,\ell}=P_i^{(\ell)}\).  Injectivity of
\(\mathbf Z[T]\to\Ql[T]\) gives \(\widetilde P_i=P_i\).  This proves
existence, factorization, and uniqueness in \textup{(ii)}.

\smallskip
\noindent\emph{Clause \textup{(iii)}.}
Fix \(0\leq i\leq2d\).  Let \(\alpha\) be a reciprocal root of \(P_i\),
and put
\(r=\deg P_i\).  Since \(P_i(0)=1\),
\[
 P_i^\#(U)=U^rP_i(U^{-1})\in\mathbf Z[U]
\]
is monic and has \(\alpha\) as a root; hence \(\alpha\) is an algebraic
integer.  By Lemma~\ref{lem:embedding-extension}, choose a
\(\mathbf Q\)-isomorphism
\(\eta:\overline{\mathbf Q}\xrightarrow{\sim}
\overline{\mathbf Q}^{(\ell)}\subset\Qlb\).  Clause
\textup{(ii)} makes \(\eta(\alpha)\) a reciprocal root of
\(P_{i,\ell}\), hence an eigenvalue of \(F_{i,\ell}\).

Let \(\tau:\mathbf Q(\alpha)\hookrightarrow\mathbf C\).  Lemma
~\ref{lem:embedding-extension}, applied to
\(\eta(\mathbf Q(\alpha))\) and \(\tau\circ\eta^{-1}\), supplies an
isomorphism \(\iota:\Qlb\xrightarrow{\sim}\mathbf C\) with
\(\iota(\eta(\alpha))=\tau(\alpha)\).  Proposition
~\ref{prop:weight-purity} gives
\[
 |\tau(\alpha)|=|\iota(\eta(\alpha))|=q^{i/2}.
\]
Thus \(\alpha\) is a \(q\)-Weil number of weight \(i\).  Let \(\beta\) be
an ordinary root of \(P_i\), and let
\(\tau:\mathbf Q(\beta)\hookrightarrow\mathbf C\).  Apply the preceding
reciprocal-root conclusion to \(\alpha=\beta^{-1}\); since
\(\mathbf Q(\alpha)=\mathbf Q(\beta)\), it gives
\(|\tau(\beta)|=q^{-i/2}\).  This proves \textup{(iii)}.

\smallskip
\noindent\emph{Clause \textup{(iv)}.}
Fix \(0\leq i\leq2d\).
Scalar extension preserves the degree of the nonzero integral polynomial
\(P_i\).  Since \(F_{i,\ell}\) is invertible, the leading coefficient of
\(P_{i,\ell}\) is nonzero, and therefore
\[
 b_i=\deg P_i=\deg P_{i,\ell}
 =\dim_{\Ql}H^i(X,\Ql).
\]
The left side is independent of \(\ell\).  Proposition
~\ref{prop:poincare-pairing} gives the Frobenius-equivariant isomorphism
\[
 H^{2d-i}(X,\Ql)\xrightarrow{\sim}H^i(X,\Ql)^\vee(-d).
\]
Taking dimensions and using the preceding degree formula in degrees \(i\) and
\(2d-i\) gives \(b_i=b_{2d-i}\).  Proposition~\ref{prop:endpoints} gives
\[
 b_0=b_{2d}=1,
 \qquad P_0(T)=1-T,
 \qquad P_{2d}(T)=1-q^dT.
\]
This proves \textup{(iv)}.

\smallskip
\noindent\emph{Clause \textup{(v)}.}
Use Definition~\ref{def:determinant-constants} to write
\[
 D_i=(-1)^{b_i}[T^{b_i}]P_i(T),
 \qquad
 \chi=\sum_{i=0}^{2d}(-1)^ib_i,
 \qquad
 \Delta=\prod_{i=0}^{2d}D_i^{(-1)^i}.
\]
Proposition~\ref{prop:poincare-pairing} supplies a perfect pairing
\[
 H^i(X,\Ql)\times H^{2d-i}(X,\Ql)\longrightarrow\Ql(-d)
\]
for which \(\langle Fx,Fy\rangle=q^d\langle x,y\rangle\).  Thus an
eigenvalue \(\alpha\) in degree \(i\) corresponds, including multiplicity,
to \(q^d/\alpha\) in degree \(2d-i\).  In middle degree the pairing is
\((-1)^d\)-symmetric; if \(d\) is odd it is a nondegenerate alternating
pairing, so \(b_d\) is even.

For every \(0\leq i\leq2d\), Proposition
~\ref{prop:duality-reciprocity} translates this eigenvalue relation to the
canonical integral factors and gives
\begin{equation}
 P_{2d-i}(q^{-d}T^{-1})
 =(-1)^{b_i}T^{-b_i}D_i^{-1}P_i(T),
 \qquad
 D_iD_{2d-i}=q^{db_i}.
 \label{eq:weil-final-reciprocity-data}
\end{equation}
Using the middle-degree parity supplied by the Poincare pairing, we have
\[
 \chi=2\sum_{0\leq i<d}(-1)^ib_i+(-1)^db_d,
\]
the integer \(d\chi\) is even: this is immediate if \(d\) is even, and
if \(d\) is odd, write \(b_d=2m\); then
\(\chi=2(\sum_{0\leq i<d}(-1)^ib_i-m)\), so \(d\chi\) is again even.
Thus \(h=d\chi/2\in\mathbf Z\).

Reindexing \(i\mapsto2d-i\) and using
Equation~\ref{eq:weil-final-reciprocity-data} gives
\[
 \Delta^2
 =\prod_{i=0}^{2d}(D_iD_{2d-i})^{(-1)^i}
 =\prod_{i=0}^{2d}q^{d(-1)^ib_i}
 =q^{d\chi}.
\]
Definition~\ref{def:functional-equation-sign} sets
\(\varepsilon=(-1)^\chi q^{-h}\Delta\).  Hence
\(\varepsilon^2=q^{-2h}\Delta^2=1\), so
\(\varepsilon\in\{1,-1\}\).

Finally, use the factorization from \textup{(ii)}, reindex, and apply the
first identity of Equation~\ref{eq:weil-final-reciprocity-data}:
\[
\begin{aligned}
 Z(X_0,q^{-d}T^{-1})
 &=\prod_{i=0}^{2d}
   P_{2d-i}(q^{-d}T^{-1})^{(-1)^{i+1}}\\
 &=\prod_{i=0}^{2d}
   \bigl((-1)^{b_i}T^{-b_i}D_i^{-1}P_i(T)\bigr)^{(-1)^{i+1}}\\
 &=(-1)^\chi T^\chi\Delta
   \prod_{i=0}^{2d}P_i(T)^{(-1)^{i+1}}\\
 &=\varepsilon q^{d\chi/2}T^\chi Z(X_0,T).
\end{aligned}
\]
In the third line we used
\(\sum_i b_i\equiv\sum_i(-1)^ib_i=\chi\pmod2\).  This proves
\textup{(v)} in \(\mathbf Q(T)\).

\smallskip
\noindent\emph{Clause \textup{(vi)}.}
Let \(\mathscr L\) be an arbitrary smooth proper lift of \(X_0\) with
associated complex fiber \(X_{\mathbf C}(\mathscr L)\), and fix
\(0\leq i\leq2d\).  Proposition
~\ref{prop:betti-comparison} gives
\[
 H^i(X,\Ql)\cong
 H^i_{\mathrm{sing}}
 \bigl(X_{\mathbf C}(\mathscr L)(\mathbf C),\mathbf Q\bigr)
 \otimes_{\mathbf Q}\Ql.
\]
Taking dimensions and using \textup{(iv)} gives
\[
 b_i=\dim_{\mathbf Q}H^i_{\mathrm{sing}}
 \bigl(X_{\mathbf C}(\mathscr L)(\mathbf C),\mathbf Q\bigr).
\]
Because \(\mathscr L\) and \(i\) were arbitrary, this proves \textup{(vi)}.  No
lifting hypothesis was used in \textup{(i)}--\textup{(v)}.
\end{proof}
